% concatenated from main_submission_3.tex
\documentclass{amsart}


\usepackage{comment}
\usepackage{xcolor}
\usepackage{amssymb}
\usepackage{graphicx}

\newcommand{\rev}[1]{\textcolor{blue}{#1}}
\newenvironment{revision}{\color{blue}}{\color{black}}
\newcommand{\jemsrev}[1]{\textcolor{red}{#1}}


\newcommand{\ri}{\mathrm{i}}
\newcommand{\eig}{\mathrm{eig}}
\newcommand{\spec}{\mathrm{spec}}
\newcommand{\sign}{\mathrm{sign}}
\newcommand{\depth}{\mathrm{depth}}
\newcommand{\diag}{\mathrm{diag}}


\theoremstyle{plain}
\newtheorem{theorem}{Theorem}[section]
\newtheorem{proposition}[theorem]{Proposition}
\newtheorem{lemma}[theorem]{Lemma}
\newtheorem{corollary}[theorem]{Corollary}
\newtheorem{assumption}{Assumption}
\theoremstyle{definition}
\newtheorem{definition}[theorem]{Definition} 
\newtheorem{example}[theorem]{Example}
\newtheorem{xca}[theorem]{Exercise}
\newtheorem{condition}[theorem]{Condition}
\newtheorem{remark}[theorem]{Remark}

\begin{document}

\title[Resonant Spectral Submanifolds]{Existence and Regularity of Stable Resonant Spectral Submanifolds and Linearization Maps}
\author{Florian Kogelbauer}
\thanks{Corresponding Author: Florian Kogelbauer}
\address{Department of Mathematics, ETH Z\"{u}rich, R\"{a}mistrasse 101, 8092 Z\"{u}rich, Switzerland}
\email{floriank@ethz.ch}
\author{Rafael de la Llave}
\address{School of Mathematics, Georgia Institute of Technology, 30332-0160 Atlanta, Georgia}
\email{rafael.delallave@math.gatech.edu}

\keywords{Resonances, Spectral Submanifold, Linear Conjugacy, Logarithmic Polynomial, Polyhomogeneous Functions, Hartman Conjecture}

\subjclass[2020]{37D10, 37C05, 37C30}

\begin{abstract}
We prove the existence of $C^{r,1-}$-regular stable invariant manifolds and linear conjugacies for analytic maps near a hyperbolic fixed point in the presence of resonances. The regularity exponent $r$ depends on the minimal resonant index, while the H\"older exponent can be chosen arbitrarily close to $1$, i.e., $1-\varepsilon$ for any $\varepsilon>0$. As a consequence, we obtain the stable version of the Hartman conjecture for analytic systems with semisimple linearization: in the fully stable case the local conjugacy can be chosen $C^{1,1-\varepsilon}$ for every $\varepsilon>0$.\\
Our approach introduces a new class of functional expansions based on logarithmic polynomials, which enables the invariance equation to be solved explicitly and algorithmically to arbitrary order. The existence results are obtained via a fixed-point argument in a suitable Banach space. We further present several analytic examples that both illustrate the applicability of the theorem while demonstrating the necessity of its assumptions.
\end{abstract}

\maketitle

\section{Introduction}


Understanding the dynamics near an equilibrium is a classical problem in dynamical systems. Given a map or vector field with a fixed point, two closely related questions arise. The first is the regularity of a local change of variables that conjugates the nonlinear dynamics to its linear part. In the hyperbolic case, the Hartman--Grobman theorem gives a topological conjugacy \cite{grobman1959homeomorphism,hartman1960lemma,hartman1960local}, which prompts the more delicate question of how regular such a conjugacy can be. The second question is the existence and regularity of invariant manifolds tangent to invariant spectral subspaces of the linearization. In the parameterization method, these two questions can be formulated through closely related invariance equations.\\

More precisely, let $f:\mathbb{C}^s\to\mathbb{C}^s$ be a sufficiently regular map with $f(0)=0$, $A=Df(0)$ and let $X$ be an $A$-invariant subspace. In the parameterization method one seeks a parametrization $K$ and a reduced dynamics $R$ satisfying $f\circ K=K\circ R$. In this paper we study the more restrictive invariance-linearization equation obtained by prescribing the reduced dynamics to be the linear one, $R=A|_X$, so that
\begin{equation}\label{invariance_eq_intro}
        f\circ K=K\circ A|_X.
\end{equation}
We require $K$ to be tangent at the origin to the inclusion of $X$ into the ambient space. Thus, a solution simultaneously parametrizes an invariant manifold tangent to $X$ and linearizes the dynamics restricted to that manifold. Its regularity immediately implies the same regularity for the corresponding invariant manifold, but the converse need not hold: the image manifold may be smoother than a parametrization solving the invariance-linearization equation, because a loss of regularity may come from the additional requirement that the dynamics on the manifold be conjugated to its linearization rather than from the geometry of the manifold itself. This distinction is important below and is also the reason that it is useful to discuss invariant manifolds and the Hartman linearization problem together.\\


The purpose of the present paper is to understand the invariance--linearization problem \eqref{invariance_eq_intro} in the presence of resonances, in the sense of Definition \ref{def_resonances}. Rather than excluding the resonant orders at which the usual cohomological equations become singular, we enlarge the Taylor ansatz by logarithmic-polynomial terms and resolve the resulting equations recursively. The logarithmic approximation is then corrected to an exact solution by a fixed-point argument in the weighted logarithmic-series Banach space $\Gamma_N^d$, introduced in Section \ref{notation}. A realization estimate, separated from the coefficient-space fixed-point argument, gives convergence in $C^{r_*,1-}$. The exponent $r_*$ is a spectrum-dependent lower guarantee for the invariance-linearization parametrization, not in general the optimal regularity of its image. If the first relevant resonant forcing coefficient vanishes, the associated logarithmic term disappears and higher regularity may occur. In the full stable case the construction yields a $C^{1,1-\varepsilon}$ linearizing conjugacy for every $\varepsilon>0$. Classical results of Hartman and subsequent refinements guarantee differentiable or $C^{1,\beta}$ linearization under contraction or spectral hypotheses, with $\beta$ tied to the regularity and spectral condition, see, in particular, \cite{hartman1960local,newhouse2017hartman,zhang2017differentiability}. Our conclusion is a near-Lipschitz H\"older estimate for the derivative under the stronger analytic and semisimplicity hypotheses used here.\\

\subsection{Some Remarks on Previous Literature}\label{intro_literature}

The role of resonances is already visible in the classical linearization theorems. Sternberg and Chen showed that, under suitable non-resonance assumptions, a sufficiently smooth hyperbolic system is smoothly locally conjugate to its linear part \cite{chen1963equivalence,chen1965local,sternberg1957local}, see also the geometric formulation in \cite{banyaga1996cohomology}. In the analytic category, small divisors lead to arithmetic conditions of Siegel--Bruno type \cite{siegel1942iteration,bruno1971analytic}. With resonances, one generally expects a normal form rather than complete analytic linearization, since resonant terms cannot in general be removed by a near-identity change of variables \cite{arnold2012geometrical,carr2012applications}. In the full-dimensional hyperbolic problem, Hartman's differentiability results led to what is often called the \emph{Hartman conjecture} \cite{hartman1960lemma,hartman1960local,pugh1969theorem}. Related developments include bunching and invariant-foliation regularity \cite{hasselblatt1994periodic,pugh1997holder} and differentiable linearization results of Belitskii, Samovol, Stowe, Van Strien, Newhouse, and others \cite{belitskii1978equivalence,samovol2010polynomial,samovol2010new,stowe1988linearization,vanstrien1990smooth,newhouse2017hartman,zhang2017differentiability}. Resonant hyperbolic linearization has also been studied by Rayskin and by Bonckaert and Naudot \cite{rayskin1998alpha,bonckaert2003linearization}.\\

For invariant manifolds, the same cohomological mechanism appears, but the regularity of the invariant image must be distinguished from the regularity of a parametrization that also linearizes its internal dynamics. For non-resonant spectral subspaces, the parameterization method gives strong existence, regularity, and computational results \cite{Cab2003pam,Cab2003,CABRE2005444}, with related developments for difference equations and numerical implementations \cite{delaLlaveLomeli2012,gonzalez2022finite}. In the stable case, the spectral subspace associated with the least contracting eigenvalues gives the slow directions governing asymptotic approach to the fixed point \cite{Wayne97}. Several notions of slow invariant manifold coexist \cite{lorenz1986existence,lorenz1987nonexistence,lorenz1992slow}, and a precise theory under suitable non-resonance assumptions was developed in \cite{Llave97,Cab2003,CABRE2005444}. Internal resonances may be absorbed into a nonlinear reduced dynamics when only the invariant manifold is sought, whereas external resonances couple selected and transverse directions and can impose genuine regularity obstructions \cite{belitskii1978equivalence,samovol2010polynomial,samovol2010new,homburg2006invariant}. Our results treat the corresponding invariance-linearization problem without excluding these external resonances.\\

Logarithmic expansions themselves, sometimes called $\Psi$-series, have a long history and appear in the work of Horn and in the classical theory of ordinary differential equations \cite{horn1896ueber,hille1997ordinary,gavrilov1992non}. Related concepts occur in invariant-manifold and normal-form problems, for instance in work of Chaperon \cite{chaperon1986c,CHAPERON_2004}, and products of polynomials and logarithms appear in \cite{Mourtada91,bonckaert2003linearization}. In geometric singular analysis, mixed power-logarithmic expansions are known as polyhomogeneous expansions \cite{grieser2001basics}. The present paper brings this viewpoint to resonant invariance equations near hyperbolic fixed points and, crucially, combines it with a Banach-space convergence argument and a recursive computational scheme.\\

There is also a direct connection with resonant forcing in mechanical systems \cite{nayfeh2024nonlinear}. In Lyapunov's expansions near an equilibrium \cite{liapounoff1907probleme}, substitution of lower-order terms produces forcing whose exponential time dependence is determined by products of eigenmodes. When such a product has the same exponential rate as another eigenmode, the inhomogeneous linear equation is resonantly forced and acquires a polynomial factor in time. In spatial coordinates this mechanism produces the logarithmic terms used in the subsequent analysis and can be regarded as the flow analogue of the singular cohomological equations at resonant multi-indices.\\

The contribution of the present paper can therefore be summarized as follows. We allow external resonances in the invariance-linearization equation and resolve each singular cohomological equation by explicit logarithmic monomials. We prove convergence to an exact solution in the Banach space $\Gamma_N^d$, whose precise definition is given in Section \ref{notation}, and we derive the resonance-dependent regularity $C^{r_*,1-}$, including finer directional regularity. The recursive construction is algorithmic to arbitrary finite order and is directly amenable to symbolic power-series manipulation. We have carried it out by hand in several examples, while a systematic software implementation is left for future work. The logarithmic class is substantially more structured than a generic $C^{r,\alpha}$ class and is smooth away from the coordinate singularities responsible for the logarithms.\\
The restriction to linear reduced dynamics is a deliberate choice of our approach. Internal resonances can often be absorbed into a nonlinear reduced map if one seeks only an invariant manifold, but composition with a prescribed nonlinear reduced dynamics is not naturally compatible with the logarithmic Banach spaces used here. A finite-order logarithmic ansatz followed by a $C^r$ argument is possible in principle, but then one must balance smallness of the remainder against the regularity of the approximate solution. We leave this extension for future work. Likewise, the logarithmic class should not be identified with a generic $C^{r,\alpha}$ class: it carries additional algebraic information about the resonant variables and admits regularity properties not encoded by the single exponent $r$.\\

We use throughout the resonance terminology and the set $I_{\rm res}(\Lambda)$ from Definition \ref{def_resonances} in Section \ref{notation}. The main result is the following.


\begin{theorem}\label{mainthm}
Let $f:\mathbb{C}^s \to\mathbb{C}^s$ be an analytic map with $f(0)=0$, understood, when $f$ represents a real system, in complex eigencoordinates compatible with the underlying real structure. All differentiability assertions below are in the real sense. Let
\begin{equation}
A=Df(0),
\end{equation}
be semisimple and invertible.
Let $X_1$ be an $A$-invariant, $d$-dimensional subspace spanned by eigenvectors of $A$ associated to eigenvalues $\Lambda = \{\lambda_1,...,\lambda_d\}$ such that
\begin{equation}\label{asslambda}
|\lambda_j|< 1,\quad 1\leq j\leq d.
\end{equation}
Assume first that $I_{\rm res}(\Lambda)\neq\varnothing$. For each resonant multi-index $n\in I_{\rm res}(\Lambda)$, define
\begin{equation}\label{def_r_star}
    r(n):=
\min\left\{
|n|-1,\,
\min_{j\in\operatorname{supp}(n)} n_j
\right\} \geq 1,\quad r_*:=
\min_{n\in I_{\rm res}(\Lambda)} r(n).
\end{equation}
Then, on a sufficiently small neighborhood of the origin in $X_1$, there exists a parametrization $h$ of class $C^{r_*,1-\varepsilon}$ for every $\varepsilon\in(0,1)$ satisfying
\begin{equation}
        h(0)=0,\qquad Dh(0)=\iota_{X_1},\qquad f\circ h=h\circ A|_{X_1},
\end{equation}
where $\iota_{X_1}$ denotes the canonical inclusion embedding of $X_1$ into the ambient space. After restricting the domain, $h$ is an embedding and its image is an $f$-invariant manifold tangent to $X_1$. The parametrization admits logarithmic-polynomial approximations to arbitrary finite order. The stated regularity is a uniform lower bound on the regularity of this invariance-linearization parametrization. If $I_{\mathrm{res}}(\Lambda)=\varnothing$, the parametrization can be chosen analytic.
\end{theorem}

\begin{remark}
Theorem \ref{mainthm} is formulated for the strictly stable part of the spectrum as expressed in assumption \eqref{asslambda}. Of course, the same conclusions hold true if we consider the strictly unstable part of the spectrum and  $f^{-1}$ instead of $f$, i.e., assumption \eqref{asslambda} can be reversed.\\
\end{remark}

\begin{remark}
The semisimplicity assumption is used in this paper to keep the spectral decomposition and the logarithmic cohomology equations transparent. The same strategy is expected to extend to general invertible $A$ after replacing eigendirections by generalized eigenspaces and keeping track of the additional nilpotent terms. We do not use that extension in the present proof and therefore state the theorem only in the semisimple setting.
\end{remark}


As explained above, the full-dimensional choice $X_1=\mathbb{C}^s$ turns the invariance-linearization equation into the usual local linearization problem. Theorem \ref{mainthm} therefore yields the following stable Hartman statement under the hypotheses used here.


\begin{corollary}\label{cor_Hartman}
    Under the assumptions of Theorem \ref{mainthm}, if $X_1$ is the full strictly stable space, then $f$ is locally $C^{1,1-\varepsilon}$-conjugate to its linearization for every $\varepsilon\in(0,1)$.
\end{corollary}

Corollary \ref{cor_Hartman} is a stable, analytic, semisimple differentiable-linearization statement. It is useful to distinguish it from the various formulations historically grouped under the Hartman conjecture. Hartman's contraction theorem gives differentiable linearization and, in finite dimensions, $C^{1,\beta}$ regularity for some $\beta>0$, while later work treats weaker smoothness and more general hyperbolic spectral configurations \cite{hartman1960local,newhouse2017hartman,zhang2017differentiability}. Under our stronger analytic and semisimplicity assumptions, the logarithmic construction gives the same conjugacy in $C^{1,\alpha}$ for every $\alpha<1$. We do not claim here a corresponding result for mixed stable--unstable spectra.\\

Our proof has two steps. First, we solve the cohomological equations recursively and construct logarithmic-polynomial approximate solutions to arbitrarily high order. Second, after all resonant orders have been incorporated, we correct the approximation by a contraction on a closed ball of the tail space $\Gamma_N^d$. The coefficient-space topology supplies the algebra and composition estimates needed for the functional equation, while the explicit power-logarithmic form yields the $C^{r_*,1-}$ realization and the finer directional regularity.\\

The paper is organized as follows. Section \ref{notation} collects notation and introduces logarithmic polynomials and the Banach spaces used in the proof. Section \ref{prelim} formulates the invariance and linearization problem and discusses the role of outer resonances. Section \ref{approxsol} constructs approximate logarithmic-polynomial solutions recursively, while Section \ref{existence} proves the existence and regularity theorem by the fixed-point argument. Section \ref{examples} gives examples illustrating the hypotheses and sharpness, while Section \ref{conclusion} contains concluding remarks and extensions.



\section{Notation and Basic Definitions}\label{notation}

Throughout the paper, $s$ will be the dimension of the ambient space, while $d\leq s$ will be the dimension of an invariant subspace or manifold. We denote the $n^{th}$ unit vector in the vector space $\mathbb{C}^s$, e.g., the vector with all zero entries and one at position $n$, as $1@n$. For a multi-index $n\in\mathbb{N}^d$, $n=(n_1,..,n_d)$, let
%we define the\emph{depth of $n$} as\begin{equation}\depth(n)=\max_{1\leq i\leq d} n_i,\end{equation} and let
\begin{equation}
    |n| = \sum_{j=1}^d n_j,
\end{equation}
denote its order and let
\begin{equation}
    \operatorname{supp}(n) = \{j: n_j\neq 0\},
\end{equation}
denote its support.\\
For a vector $x=(x_1,...,x_d)$, we denote the $(d-1)$-dimensional vector whose $x_k$-component has been removed as $\underline{x}_k$. We denote the $d\times d$ diagonal matrix with elements $a_1,...,a_d$ as $\diag(a_1,...,a_d)$. The ball of radius $r$ in $\mathbb{R}^d$ is denoted as $B^r({\mathbb{R}^d})$, while the unit ball in $\mathbb{C}^d$ is denoted as $B^r({\mathbb{C}^d})$. When the distinction between real and complex vector spaces are clear from the context, we simply write $B^r_d$.\\
We denote the $M$-dimensional complex strip of width $\sigma$ as 
\begin{equation}\label{def_strip}
    \mathfrak{T}_\sigma = \{z\in \mathbb{C}^M: |\Im(z_j)|\leq \sigma,\quad 1\leq j \leq M \}. 
 \end{equation}
%For a complex function $H:\mathbb{C}^d\to\mathbb{C}^s$, we denote the real derivative, regarding $H$ as a function from $\mathbb{R}^{2d}$ to $\mathbb{R}^{2s}$, as $D_\mathbb{R}H$.
Let $\alpha\in (0,1]$. A function $f:U\to\mathbb{R}^{d_2}$ defined on an open subset $U\subseteq\mathbb{R}^{d_1}$ is $\alpha$-H\"{o}lder continuous if
\begin{equation}
    |f(x)-f(y)|\leq C_U |x-y|^\alpha,\quad x,y\in U,
\end{equation}
for some constant $C_U$ depending on the domain. If $\alpha=1$, the function is Lipschitz continuous. A function $f\in C^{k,\alpha}(U,\mathbb{R}^{d_2})$ is $k$-times differentiable with $k^{th}$ derivative $\alpha$-H\"{o}lder continuous. We write $C^{r,1-}$ for the class of functions $C^{r,1-\varepsilon}$ for all $\varepsilon\in (0,1)$, i.e.,
\begin{equation}\label{def_C_1_minus}
    C^{r,1-} = \bigcap_{\varepsilon\in (0,1)} C^{r,1-\varepsilon}.
\end{equation}
While each $C^{r,\alpha}$ is a Banach space under the H\"{o}lder norm 
\begin{equation}\label{def_C_r_norm}
\|f\|_{C^{r,\alpha}(\Omega)}
=
\sum_{|\beta|\leq r}\|D^\beta f\|_{L^\infty(\Omega)}
+
\sum_{|\beta|=r}
\sup_{\substack{x,y\in\Omega\\ x\neq y}}
\frac{|D^\beta f(x)-D^\beta f(y)|}{|x-y|^\alpha},
\end{equation}
where $f:\Omega\to\mathbb{R}^d$ for $\Omega\subseteq\mathbb{R}^d$ open and connected, the space \eqref{def_C_1_minus} is only a Fr\'{e}chet space.\\
Let $\mathcal{P}_{d,K}$ denote the space of polynomials in $d$ variables of degree at most $K$ for which  
\begin{equation}
\dim \mathcal{P}_{d,K} = \binom{K+d}{d}.
\end{equation}
We write
\begin{equation}
\mathcal{P}_d = \bigcup_{K = 0}^\infty  \mathcal{P}_{d,K}
\end{equation}
for the vector space of polynomials in $d$ variables.  \\
A \emph{power function} is a generalized polynomial expression $\mathfrak{p}:\mathbb{C}^d\to \mathbb{R}^+$ of the form\footnote{We choose to denote power functions with Gothic fonts to cleanly distinguish them from polynomial expressions and avoid potential confusions.}
\begin{equation}\label{defpower}
    \mathfrak{p}(x) = \sum_{j=1}^d p_j  |x_j|^{\nu_j},
\end{equation}
where $\nu_j>0$ and $p_j>0$. If all the $\nu_j$ are non-negative even integers, power functions reduce to a particular class of ordinary polynomials.\\
For an operator $T:X\to Y$ defined on Banach spaces $(X,\|.\|_X)$ and $(Y,\|.\|_Y)$, we write
\begin{equation}
    \|T\|_{\rm op} = \sup_{\|x\|_{X}=1} \|Tx\|_{Y},
\end{equation}
for its operator norm and 
\begin{equation}
    \sigma(T) = \{\lambda \in \mathbb{C}: T-\lambda \text{ Id} \text{ is not bijective with a bounded inverse}\},
\end{equation}
for its operator spectrum. 



\begin{definition}\label{def_resonances}
Given an $s\times s$ matrix $A$ over the complex numbers, let $\sigma(A)\subset\mathbb{C}$ denote its spectrum and $\lambda=(\lambda_1,\ldots,\lambda_s)$ its eigenvalues, counted with multiplicity. We call a multi-index $n\in\mathbb{N}^s$ with $|n|\geq 2$ \emph{resonant} if there exists an eigenvalue $\widetilde{\lambda}$ of $A$ such that
\begin{equation}\label{lambdares}
    \lambda^n=\widetilde{\lambda}.
\end{equation}
If a multi-index is not resonant, we call it \emph{non-resonant}.\\

Let $\Lambda=\{\lambda_1,\ldots,\lambda_d\}\subseteq\sigma(A)$ be a subset of the spectrum of $A$. For a multi-index $n\in\mathbb{N}^d$ with $|n|\geq 2$, we write
\begin{equation}
    \lambda_\Lambda^n:=\lambda_1^{n_1}\cdots\lambda_d^{n_d}.
\end{equation}
We call $n$ \emph{resonant with respect to $\Lambda$} if there exists $1\leq k\leq s$ such that
\begin{equation}\label{lambdaresLambda}
    \lambda_\Lambda^n=\lambda_k.
\end{equation}
Such a resonance is called an \emph{internal resonance} or \emph{inner resonance} relative to $\Lambda$ if $\lambda_k\in\Lambda$, otherwise it is called an \emph{external resonance} or \emph{outer resonance}. For a given subset $\Lambda=\{\lambda_1,\ldots,\lambda_d\}\subseteq\sigma(A)$, we denote the set of all its resonant indices by
\begin{equation}\label{def_I_res}
    I_{\rm res}(\Lambda)
    =
    \left\{
        n\in\mathbb{N}^d :
        |n|\geq2,\quad
        \lambda_\Lambda^n=\lambda_k
        \text{ for some } 1\leq k\leq s
    \right\}.
\end{equation}
\end{definition}

\begin{lemma}\label{lemma_finite_resonances}
Under \eqref{asslambda} and the invertibility of $A$, the set $I_{\rm res}(\Lambda)$ is finite.
\end{lemma}
\begin{proof}
Set $\rho=\max_{1\leq j\leq d}|\lambda_j|<1$ and $m=\min_{\lambda\in\sigma(A)}|\lambda|>0$. If $n\in I_{\rm res}(\Lambda)$, then $m\leq|\lambda_\Lambda^n|\leq\rho^{|n|}$. Hence $|n|\leq (\log m)/(\log\rho)$, with the inequality interpreted in the usual way since $\log\rho<0$. Thus only finitely many multi-indices can be resonant.
\end{proof}


\begin{remark}
In the following, the results are written in complex eigencoordinates for notational convenience and are intended for real systems after complexification: the real linear part $A$ has non-real eigenvalues in conjugate pairs, and all differentiability assertions are understood in the real sense. For such systems, one can always re-write the complex Jordan normal form by using two-by-two real block matrices of the form
\begin{equation}
    \begin{pmatrix} \Re(z)& -\Im(z) \\  \Im(z) & \Re(z)  \end{pmatrix} ,
\end{equation}
for the complex variable $z$. This is equivalent to regarding the complex space $\mathbb{C}^d$ as $\mathbb{R}^{2d}$. The use of the complex notation simplifies the typography of 
many algebraic formulas. On the other hand, bearing in mind that complex derivatives are different as compared to real derivatives, we will always understand the regularity of functions in the real sense.
\end{remark}

\subsection{Logarithmic Series}

Before we state the problem of invariance, conjugacy and linearization, we introduce some classes of functions and appropriate Banach spaces associated with them. An expression of the form
\begin{equation}\label{deflogpoly}
\sum_{1\leq |n|\leq N} H_n(\log|\phi_1(x)|,...,\log|\phi_M(x)|)x^n,
\end{equation}
where $\phi_1,...,\phi_M$ are H\"{o}lder continuous functions in $x$ in a neighborhood of the origin and $H_n$ are polynomial functions is called \emph{logarithmic polynomial of degree $N$}. For $N=\infty$, an expression of the form \eqref{deflogpoly} is called \emph{logarithmic series}. The functions $\phi_j$ are called \emph{basis functions}, while $\eta_j:=\log|\phi_j|$ will be called the corresponding \emph{logarithmic variables}.

\begin{remark}
    Typically, it is assumed that the functions $\phi_j$ are at most of power-law growth at infinity,
\begin{equation}\label{polybound}
    \phi(x) = \mathcal{O}(|x|^\nu),\quad x\to\infty,
\end{equation}
for some $\nu>0$. In the following, however, we will be mostly interested in the behavior of functions of the form \eqref{deflogpoly} around the origin, so we do not include the asymptotic bound \eqref{polybound} in the definition of a logarithmic series. We stress, on the other hand, that the functions $\mathfrak{p}_j$ we are using in the explicit construction of approximate solutions to the invariance equation in the presence of resonances will be power functions of the form  \eqref{defpower}, which, of course, satisfy \eqref{polybound}. Let us also remark the assumption of $\phi_j$ being H\"{o}lder continuous is not substantial for the subsequent arguments. Again, using power functions for $\phi_j$ of course implies a certain regularity, which is at least H\"{o}lder continuous.  
 \end{remark}


For the fixed-point argument appearing later in the paper, we will need appropriate function spaces that reflect the regularity of solutions as well as proximity to approximate solutions to the invariance equation, which will be given in terms of logarithmic polynomials.\\
For any $\beta,\sigma >0$, we define the following family of norms on the space of analytic functions\footnote{For the following definitions and arguments, it is actually irrelevant if we define the norm \eqref{defbetanorm} specifically on analytic functions,i.e., the completion of polynomials under a suitably chosen norm, since it will act on polynomials of a maximal degree exclusively.} $f:\mathbb{R}^d\to \mathbb{C}^s$, 
\begin{equation}\label{defbetanorm}
\|f\|_{\beta,\sigma} = \sup_{y\in \mathfrak{T}_\sigma} e^{-\beta |\Re y|}|f(y)|,
\end{equation}
see \eqref{def_strip}. Further, consider the vector space of series of the form 
\begin{equation}\label{def_H}
H = \sum_{|n|=1}^\infty H_n(\log|\phi_1|,...,\log|\phi_M|)x^n,\quad H_n:\mathfrak{T}_\sigma\to\mathbb{C}^s \text{ analytic },
\end{equation}
defined on the strip $\mathfrak{T}_\sigma\subset \mathbb{C}^M$ and introduce the following $l^1$-type norm
\begin{equation}\label{defnorm}
    \|H\|_{\Gamma_N^d,\beta,\sigma} = \sum_{|n| = N}^{\infty} \|H_n\|_{|n|\beta,\sigma}. 
\end{equation}

The space 
\begin{equation}
 C^\omega_{\beta,\sigma}(\mathbb C^M)
    :=
    \left\{
    f\text{ analytic in }  \mathfrak{T}_\sigma:
   \|f\|_{\beta,\sigma}<\infty
    \right\},
\end{equation}
is a Banach space and thus the space 
\begin{equation}\label{defGamma}
    \Gamma^d = \left\{ H \text{ is of the form \eqref{def_H} and }  \|H\|_{\Gamma_N^d,\beta,\sigma} <\infty \right\},
\end{equation}
is a Banach space as well. We denote the linear subspace of logarithmic series starting at order $N$ as 
\begin{equation}
    \Gamma^d_N = \left\{H = \sum_{|n|= N }^\infty H_n(\log|\phi_1|,...,\log|\phi_M|)x^n,\quad H \in \Gamma^d \right\},
\end{equation}
also regarded as functions defined on $\mathfrak{T}_\sigma\subset \mathbb{C}^M$.\\
The Banach space $\Gamma^d_N$ is even a Banach algebra: The norm \eqref{defnorm} is somewhat similar to the $l^1$-norm on the space of analytic functions.\\
Indeed, we first note that for two polynomials $p\in\mathcal{P}_{d,n}$, $q\in\mathcal{P}_{d,m}$ and $\beta = \beta_1+\beta_2$, we have that
\begin{equation}
\begin{split}
       \|qp\|_{\beta,\sigma} & = \sup_{y\in\mathfrak{T}_\sigma} e^{-\beta |\Re y|} |p(y)q(y)|\\
       & = \sup_{y\in\mathfrak{T}_\sigma} e^{-\beta_1 |\Re y|} |p(y)| e^{-\beta_2 |\Re y|}|q(y)|\\
    & \leq \|p\|_{\beta_1,\sigma}\|q\|_{\beta_2,\sigma}.
\end{split}
\end{equation}



%where the weight satisfies the sub-additivity property \begin{equation}    w(n+m)\leq w(n) w(m),\quad n,m\in\mathbb{N}^d,\end{equation}We denote by $\Gamma^d_{N,w}$ all elements in $\Gamma_N^d$ for which \eqref{defnorm} is finite.\\

Expanding two elements $H(x)$ and $\tilde{H}(x)$ in a logarithmic series,
\begin{equation}
    H(x) = \sum_{|n|\geq N} H_{n}(x) x^n,\quad \tilde{H}(x) = \sum_{|n|\geq N} \tilde{H}_{n}(x) x^n
\end{equation}
we have that 
\begin{equation}
\begin{split}
\|H\tilde{H}\|_{\Gamma_N^d,\beta,\sigma} & = \sum_{|k|\geq 2N}\left\|\sum_{n+m=k} H_{n}\tilde{H}_m\right\|_{\beta |k|,\sigma}\\
& \leq \sum_{|k|\geq 2N}\sum_{n+m=k} \|H_{n}\tilde{H}_m\|_{\beta |k|,\sigma} \\
& \leq  \sum_{|k|\geq 2N}\sum_{n+m=k} \|H_n\|_{\beta|n|,\sigma}  \|\tilde{H}_m\|_{\beta|m|,\sigma}  \\
 & \leq  \left(\sum_{|n|\geq N} \|H_n\|_{\beta|n|,\sigma}\right) \left(\sum_{|m|\geq N} \|\tilde{H}_m\|_{\beta|m|,\sigma}\right)\\
& \leq \|H\|_{\Gamma_{N}^d,\beta,\sigma}\|\tilde{H}\|_{\Gamma_{N}^d,\beta,\sigma},
\end{split}
\end{equation}
where we have used that $k=n+m$ allows us to split $\beta|k| = \beta|n|+\beta|m|$.


\begin{remark}
    The norm \eqref{defnorm} only measures the polynomial expressions $H_n$ and does not take into account the regularity of the composed map
\begin{equation}\label{Hncomplog}
        x\mapsto H_n(\log|\phi_1(x)|,...,\log|\phi_M(x)|).
    \end{equation}
Indeed, the regularity of the function \eqref{Hncomplog} and hence of sums and series expressions of the form \eqref{deflogpoly} depends on the specific choice of functions $\phi_j$, as detailed in the following section. For the resolution of the cohomology equation in the presence of resonances, the basis functions will be of genuinely lower order compared to polynomials and thus the separation in a polynomial and a logarithmic component as measured by the norm \eqref{defnorm} is meaningful.  %We further note that a given function $f$ might admit several series expansions as written in full generality in \eqref{deflogpoly}, as shown by $x^2 = \log(e^x)x$. For the resolution of the cohomology equation in the presence of resonances, however, the basis functions will be of genuinely lower order compared to polynomials and thus the separation in a polynomial and a logarithmic component as measured by the norm \eqref{defnorm} is meaningful. 
\end{remark}

%\todo[inline, color = red]{Comment more on strangeness of the norm, choose eg $\mathfrak{p}_j(x) = e^{x_j}$, then the coefficient polynomials are actual polynomials again. We may write the function $f(x)=x^2$ as $f(x)  = \log(e^x)x $ \begin{equation}    \|f\|_{\Gamma^1,\beta} = \|1\|_{2\beta} = 1 \end{equation}$H_1(y) = y$\begin{equation}  \|f\|_{\Gamma^1,\beta} = \|H_1\|_{1\beta} = 1/\beta \end{equation}The idea is of course that the $H_n's$ are of lower order compared to $x^n$ in the following, the $H_n$ will be determined by algebraic constraints coming from the resonances, which have to be genuinely logarithmic terms of lower order. }







%\begin{remark}   Since the space of polynomials in $d$ variables of degree $n$ is finite-dimensional, every two metric on this space are equivalent, but the constants may depend on $d$ and $n$. \end{remark}







\section{Preliminaries and Formulation of the Problem of Linear Conjugacy}\label{prelim}

In this section, we precise the overall set-up for our analysis. We define the invariance-linearization equation and show the equivalence of the flow map formulation with the vector field formulation for differential equations. 



\subsection{Formulation of the Problem}

We will be concerned with analytic maps $f:\mathbb{C}^s\to\mathbb{C}^s$ such that $f(0)=0$, write
\begin{equation}\label{defA}
A=Df(0),
\end{equation}
for the linear part at zero, and set
\begin{equation}
f(x)=Ax+N_f(x),
\end{equation}
where $N_f(x)=\mathcal{O}(|x|^2)$. For a given map $f:\mathbb{C}^s\to\mathbb{C}^s$ and an $A$-invariant $d$-dimensional subspace $X_1\subseteq\mathbb{C}^s$, the general parameterization method seeks an embedding $K:B^1_d\to\mathbb{C}^s$ and a reduced map $R:B^1_d\to X_1$ satisfying $f\circ K=K\circ R$. In this paper we first impose the linear choice $R=A|_{X_1}$ and denote the corresponding parametrization by $h:B^1_d\to\mathbb{C}^s$, where $B^1_d\subset X_1$. Thus we consider the \emph{invariance-linearization equation}:
\begin{equation}\label{invariance}
f(h(x))=h(Ax),
\end{equation}
where $A=Df(0)$ and $Ax$ has to be read as $A|_{X_1}x$. Equation \eqref{invariance} states that, on the complex manifold invariant by $f$ and parametrized by $h$, the map can be conjugated to its linearization at zero. Equation \eqref{invariance} simultaneously imposes invariance of the image of $h$ and linearization of the reduced dynamics on that image and is hence called the \emph{invariance-linearization equation}.\\
%The solution $h$ defines a map analogous to a conjugacy in linearization theorems and might thus be called \emph{partial linearization map}.\\
Inner and outer resonances are understood in the sense of Definition \ref{def_resonances}. Inner resonances do not substantially affect the regularity of the invariant manifold and can be resolved with nonlinear terms in the conjugated dynamics \cite{CABRE2005444}. Indeed, allowing for more general reduced dynamics 
\begin{equation}
    f(h(x)) = h(g(x)),
\end{equation}
where 
\begin{equation}\label{nonlinearreduced}
    g(x) = A|_{X_1}x + \mathcal{O}(|x|^2),
\end{equation}
allows for nonlinear contributions and thus resolves inner resonances. Outer resonances, however, do limit the regularity of the invariant manifolds around the origin. This obstruction is crucial in certain normal form calculations \cite{Neild20140404}.\\
In the following, the distinction between inner and outer resonances will be softened. The regularity of the invariant manifold in the presence of outer resonances will be determined by a minimal resonant index and we do not make a distinction between an inner or an outer resonance, since we are interested in conjugation to the linearized dynamics. If the minimal resonant index is, indeed, an inner resonance, the regularity of the invariant manifold could be improved by considering a non-linear reduced dynamics \eqref{nonlinearreduced}.\\
Throughout the paper, as mentioned before, we use complex coordinates, mostly generalized eigendirections, even if the initial dynamical system is real. This will allow us to use Jordan canonical forms in the proofs rather than conjugate to two-by-two real block matrices, which simplifies the formulas for the approximate solutions. We stress already at this point, however, that, even though the flow maps and the approximate invariant manifolds are defined for complex arguments, they will not be complex-differentiable, i.e. analytic,  in the presence of resonances. Indeed, we will see that the complex invariant manifolds obtained in the following are only finitely differentiable in the real sense in the presence of resonances.\\
We formulate our results for maps - the relation to differential equations is well-known and treated in Appendix \ref{sec_appendix}.





\section{Approximate Solution to the Invariance Equation by Logarithmic Series}\label{approxsol}

In this section, we construct approximate solutions to the invariance equations in terms of logarithmic polynomials. These will serve as a basis for the fixed-point argument in a later section. First, we prove some results on the invertibility of certain shift operators on the space of polynomials. Then, we introduce a special choice of basis function in the logarithmic series \eqref{deflogpoly}, namely the class of power functions \eqref{defpower}, and discuss their regularity properties. Finally, we show that the invariance equation \eqref{invariance} can be solved to any order for this choice of logarithmic series.\\


\subsection{Shifted Polynomials}

Before we can construct approximate solutions to the invariance equation in terms of logarithmic series, we need some preparatory results on shifted polynomials. We emphasize already at this point that Lemma \ref{shifted} gives a \emph{unique} solution in the absence of resonances, while it only guarantees the existence of \emph{a} solution in the presence of resonances. Indeed, the lack of uniqueness in the presence of outer resonances is due to this existence of a kernel of a certain operator acting on polynomials, as elaborated in the following. While we only need invertibility properties on shifts of polynomials in one variable, we prove the general case for completeness in the following. \\

For $a\in\mathbb{C}$ and $v\in\mathbb{R}^d$, let $T_v:\mathcal{P}_d\to\mathcal{P}_d$ denote the shift operator
\begin{equation}\label{def_shift_v}
    T_v[p](x)=p(x+v),
\end{equation}
and define the operator $S_{a,v}:\mathcal{P}_d\to\mathcal{P}_d$, $S_{a,v}=I-aT_v$, i.e.,
\begin{equation}\label{defSav}
    S_{a,v}[p](x) = p(x)- a p(x+v). 
\end{equation}



%Before we proceed with a more in-depth analysis of the operator $S_{a,v}$,
We observe that the operator norm of $T_v:\mathcal{P}_{d,n}\to\mathcal{P}_{d,n}$, where $\mathcal{P}_{d,n}$ is endowed with the norm $\|.\|_{\beta,\sigma}$, is bounded as 
\begin{equation}\label{boundnormTv}
    \|T_v\|_{\rm op} \leq e^{\beta |v|},
\end{equation}
independently of the degree of $p$. Indeed, we calculate
\begin{equation}
    |p(y+v)|e^{-\beta |y|} = \Big( |p(y+v)|e^{-\beta |y+v|} \Big)e^{\beta (|y+v|-|y|)},
\end{equation}
and using that $e^{\beta (|y+v|-|y|)}\leq e^{\beta |v|}$ by the triangle inequality, we can estimate
\begin{equation}
    \begin{split}
        \|T_v p\|_{\beta,\sigma} & = \sup_{y\in\mathfrak{T}_\sigma}  |p(y+v)|e^{-\beta |y|}\\
        & \leq \Big(\sup_{y\in\mathfrak{T}_\sigma}  |p(y+v)|e^{-\beta |y+v|}\Big) e^{\beta |v|}\\
        & \leq \|p\|_{\beta,\sigma} e^{\beta|v|}
    \end{split}
\end{equation}
which immediately implies the bound \eqref{boundnormTv}. 

\begin{lemma}\label{shifted}
 Consider the operator $S_{a,v}:\mathcal{P}_{d,K}\to\mathcal{P}_{d,K}$ equipped with the norm $\|.\|_{\beta}$ as defined in \eqref{defSav} and assume that $v\neq 0$. For $a\neq 1$, the spectrum of $S_{a,v}$ and its inverse are given by 
 \begin{equation}\label{sigmaSav}
 \sigma(S_{a,v}) = \{1-a\},\quad  \sigma(S_{a,v}^{-1}) = \left\{\frac{1}{1-a}\right\},
 \end{equation}
 independently of the order of the polynomials $K$ and the shift vector $v$.\\
 In particular, $S_{a,v}:\mathcal{P}_{d,K}\to\mathcal{P}_{d,K}$ is invertible for $a\neq 1$ and we can bound its inverse as 
     \begin{equation}\label{boundSavinvfull}
    \|S_{a,v}^{-1}\|_{\rm op} \leq \frac{1}{|1-a|}\frac{(q\Delta)^{K+1}-1}{q\Delta-1},
\end{equation}
where we have set
\begin{equation}
    q = \frac{|a|}{|1-a|},\quad \Delta =\|T_v-I\|_{\rm op}.
\end{equation}
If in addition
\begin{equation}\label{boundaevb}
   |a| e^{\beta|v|}<1, 
\end{equation}
we obtain the bound
\begin{equation}\label{boundSavinv}
    \|S_{a,v}^{-1}\|_{\rm op} \leq \frac{1}{1- |a| e^{\beta|v|}},
\end{equation}
while for
\begin{equation}
|a|^{-1} e^{\beta |v|}<1,
\end{equation}
we obtain the bound
\begin{equation}\label{boundSavbinv2}
\|S_{a,v}^{-1}\|_{\rm op}\leq \frac{1}{|a|e^{-\beta|v|}-1}.
\end{equation}
For $a=1$ and $v\neq 0,$ the operator $S_{a,v}:\mathcal{P}_{d,K+1}\to\mathcal{P}_{d,K}$ is surjective. 
\end{lemma}
\begin{proof}
%Let $T_v$ denote the shift operator such that $S_{a,v}=I-a T_v$ and 
Let $D_v=v\cdot\nabla$ denote the streaming operator with velocity $v$, i.e., the generator of $T_v$. The differential operator $D_v:\mathcal{P}_{d,K}\to\mathcal{P}_{d,K}$ is well-defined and satisfies $D_v^{K+1}=0$ and hence
\begin{equation}
    T_v = \exp(D_v) = \sum_{n=0}^K \frac{D^n_v}{n!}.
\end{equation}
The operator $T_v:\mathcal{P}_{d,K}\to\mathcal{P}_{d,K}$ is thus the sum of the identity plus a nilpotent operator, and hence an unipotent operator for every $K$. This implies that $\sigma(T_v)=\{1\}$ on the whole of $\mathcal{P}_d$ and \eqref{sigmaSav} follows.\\
Writing
\begin{equation}
    S_{a,v} = (1-a)\left(I-\frac{a}{1-a}(T_v-I)\right),
\end{equation}
and noting that the operator $T_v-I$ is nilpotent of order $K+1$, i.e., $(T_v-I)^{K+1}=0$, implies that we can write a closed-form expression for the inverse of $S_{a,v}$:
\begin{equation}\label{inversexplicit}
    S_{a,v}^{-1} = \frac{1}{1-a} \sum_{k=0}^K \left(\frac{a}{1-a}\right)^k(T_v-I)^k.
\end{equation}
This immediately implies the bound 
\begin{equation}
    \|S_{a,v}^{-1}\|_{\rm op} \leq \frac{1}{|1-a|}\sum_{k=0}^K \left(\frac{|a|}{|1-a|}\right)^{k} \|T_v-I\|_{\rm op}^k,
\end{equation}
and, by summing the geometric sum, we obtain \eqref{boundSavinvfull}.\\ 
Assuming now \eqref{boundaevb}, we can invert $S_{a,v}$ by a Neumann series and the estimate \eqref{boundSavinv} follows right away. The other estimate \eqref{boundSavbinv2} follows from rewriting $S_{a,v}p=q$ as 
\begin{equation}
\frac{1}{a}p(x-v) -p(x) = \frac{1}{a}q(x-v). 
\end{equation}
To prove the surjectivity of $S_{1,v}:\mathcal{P}_{d,K+1}\to\mathcal{P}_{d,K}$, let us first note that we can factor
\begin{equation}
\begin{split}
       S_{1,v} &  = 1 - \exp(D_v)\\
       & =  1 - \sum_{n=0}^{K+1}\frac{D^n_v}{n!} \\
       & = -D_v\left(\sum_{n=0}^{K}\frac{D_v^n}{(n+1)!}\right)\\
      & =: D_v \Sigma_v,\qquad \Sigma_v:=-\sum_{n=0}^{K}\frac{D_v^n}{(n+1)!},
       \end{split}
\end{equation} 
since $D_v^{K+2}=0$ on $\mathcal{P}_{d,K+1}$. The operator $-\Sigma_v$ is the identity plus a nilpotent operator and is therefore invertible - hence $\Sigma_v$ is invertible as well. It follows that
\begin{equation}
    \text{range}(S_{1,v}) = \text{range}(D_v\Sigma_v) = \text{range}(D_v). 
\end{equation}
Surjectivity of $S_{1,v}:\mathcal{P}_{d,K+1}\to\mathcal{P}_{d,K}$ is thus equivalent to surjectivity of the streaming operator $D_v:\mathcal{P}_{d,K+1}\to\mathcal{P}_{d,K}$. The kernel of $D_v$ consists exactly of those polynomials which are constant along the $v$-direction, i.e., polynomials in the $d-1$ independent coordinates transverse to $v$. Hence
\begin{equation}
    \dim \ker D_v = \binom{K+d}{d-1}. 
\end{equation}
Therefore, the rank of $D_v$ is given by
\begin{equation}
\begin{split}
    \text{rank } D_v & = \dim \mathcal{P}_{d,K+1} - \dim \ker D_v\\
     & = \binom{K+1+d}{d} - \binom{ K+d}{ d-1}\\
     & = \binom{K+d}{d},
    \end{split}
\end{equation}
where in the last step, we have used Pascal's identity for binomial coefficients and for the dimension of the space of polynomials we refer to \cite{cox2015ideals}. Consequently,  $\text{rank } D_v = \dim \mathcal{P}_{d,K} $ and $D_v$ is surjective. 
\end{proof}


\subsection{Power Functions as Basis Functions}

%\todo[inline, color = green]{There are two classes of polynomials we have to consider: the coefficient polynomials of the logarithmic series $H_n$ which are defined by the dynamical system and the basis polynomials $p_j$ which we can choose. The form of the polynomials $H_n$ is of course relative to the basis polynomials $p_j$. We cannot choose the $p_j's$ completely freely of course, since they have to transform appropriately under the action of A in order to solve the cohomology equation. }

To resolve the invariance equation in the presence of resonances, we introduce series expansions of the form
\begin{equation}\label{expandheta}
    h(x) = \sum_{|n|=1}^\infty h_n(\eta) x^n,
\end{equation}
where each coefficient $h_n$ is a polynomial in the logarithmic variables $\eta(x)=(\eta_1(x),\ldots,\eta_M(x))$. Here $M$ is the number of logarithmic variables introduced by the resonant construction below, and
\begin{equation}\label{def_eta}
    \eta_l(x) = \log |\mathfrak{p}_l(x)|,\quad \mathfrak{p}_l(x) = \sum_{j=1}^d p_{l,j}|x_j|^{\nu_j},
    %,\quad 1 \leq j\leq M,
\end{equation}
for exponents $\nu_j>0$ and weights $p_{l,j}\geq0$ chosen below. For each $l$, at least one $p_{l,j}$ is positive; later the support of the vector $(p_{l,1},\ldots,p_{l,d})$ is chosen to match the support required by the corresponding resonant multi-index.
%We stress that the expressions of the form $|x_j|^{\nu_j}$ are interpreted as $(|x_j|^2)^{\nu_j/2}$, i.e., 
Expansion \eqref{expandheta} is a special instance of the logarithmic series \eqref{deflogpoly} for one particular choice of basis function. In order to resolve the cohomology equation in the presence of resonances, the function $x\mapsto \eta(x)$ has to satisfy the functional equation
\begin{equation}\label{func_eq_eta}
    \eta(Ax) = b + \eta(x),
\end{equation}
for some $b\neq 0$. The parameter $b$ will translate the vector quantity $v$ from the shift operator \eqref{def_shift_v}. Assuming that the coordinates $(x_1,...x_d)$ are chosen along eigendirections of $A$ without loss of generality, \eqref{func_eq_eta} implies that the exponents of the power functions necessarily have to be of the form
\begin{equation}\label{nu_form}
    \nu_j = \frac{\log(\gamma)}{\log|\lambda_j|},\quad 1\leq j\leq d,\quad 0<\gamma<1. 
\end{equation}
Indeed, we have that 
\begin{equation}
\begin{split}
    \eta(Ax) & = \log\left(\sum_{j=1}^d p_j|\lambda_j x_j|^{\nu_j}\right) =  \log\left(\sum_{j=1}^d p_j | x_j|^{\nu_j} \exp\left(\log|\lambda_j|\frac{\log(\gamma)}{\log|\lambda_j|}\right)\right)\\
    & = \log\left(\gamma\sum_{j=1}^d p_j| x_j|^{\nu_j}\right) = \log(\gamma) + \eta(x),
    \end{split}
\end{equation}
and hence, \eqref{def_eta} satisfies \eqref{func_eq_eta} for
\begin{equation}\label{defvlambda}
b = \log(\gamma),
%v_\lambda = (\log|\lambda_1|,...,\log|\lambda_d|). 
\end{equation}
as defined in \eqref{nu_form}. Subsequently, we will choose $\gamma$ adequately to guarantee certain regularity properties of $\eta$, see Lemma \ref{lemma_regularity} for details. The numerical values of the coefficients $p_j$ will be irrelevant for the further analysis, as long as $p_j>0$ for the right support index set, see Lemma \ref{lemma_regularity}. 




%\begin{remark}A real polynomial $p$ with $Z(p)=0$ has to be sign definite and we may write\begin{equation}     p(x) = \pm\sum_{j=1}^m |p_j(x)|^{2},\end{equation}for some polynomials $p_j$, see \cite{}. \end{remark}\todo[inline, color = green]{Does the same hold true for power functions?}

\begin{remark}
To solve the cohomology equation with resonances, it is the logarithmic variables $\eta_j=\log|\phi_j|$, rather than the basis functions $\phi_j$ themselves, that must satisfy an additive shift relation
\begin{equation}\label{fcteqeta}
    \eta_j(Ax)=b_j+\eta_j(x),\quad 1\leq j\leq M,
\end{equation}
for some real numbers $b_j\neq0$. Such an $\eta_j$ cannot extend continuously with a finite value to the origin: formally evaluating the relation at $x=0$ would give $\eta_j(0)=b_j+\eta_j(0)$. This is precisely why we write $\eta_j=\log|\phi_j|$ with simpler basis functions $\phi_j$ that vanish at the relevant singular set. In the construction below we take $\phi_j=\mathfrak p_j$, where $\mathfrak p_j$ is a power function of the form \eqref{defpower}. Then $\eta_j=\log\mathfrak p_j$ satisfies \eqref{fcteqeta}. This convention keeps the continuous basis functions $\phi_j$ distinct from their singular logarithmic variables $\eta_j$.

\end{remark}











%We may generalize this type of polynomials \begin{equation}    q(x) = \pm\sum_{j=1}^m |q_j(x)|^{\nu_j}.\end{equation}

%A homogeneous polynomial is a polynomial where every single term (monomial) has the same total degree, meaning the sum of the exponents for the variables in each term is identical.\\







\begin{lemma}\label{lemma_regularity}
Let $\mathfrak p:\mathbb C^d\to\mathbb R$ be a power function of the form \eqref{defpower} and let
\begin{equation}
        S=\{j\in\{1,\ldots,d\}:p_j>0\},
\end{equation}
be the support of the weights. Let $n\in\mathbb N^d$ be a multi-index with
$|n|\geq 2$ and assume the support condition $S\subseteq \operatorname{supp}(n)$, i.e.,
\begin{equation}\label{asslemma}
        p_j>0\implies n_j>0.
\end{equation}
Set
\begin{equation}\label{regularity_index}
        r
        =
        \min\left\{
        \left(\sum_{j\in S} n_j\right)-1,\,
        \min_{j\in S} n_j
        \right\},\quad |n|\geq 2,
\end{equation}
and assume moreover that
\begin{equation}\label{ass_nu_large}
        \nu_j>r+1,
        \qquad j\in S,
\end{equation}
are positive real numbers that define the exponents of the basis functions \eqref{def_eta}.
Then, for every polynomial $Q$, the function
\begin{equation}
        x\mapsto x^n Q(\log\mathfrak p(x))
\end{equation}
extends to a $C^{r,1-\varepsilon}$-function in a neighborhood of the origin, for every
$\varepsilon\in(0,1)$.
%Moreover, the extension is $C^{r+1}$ at the origin in the sense that its Taylor polynomial of order $r+1$ at the origin exists and is identically zero.
\end{lemma}

\begin{proof}
We identify $\mathbb C^d$ with $\mathbb R^{2d}$ by writing
\begin{equation}
        x_j=a_j+\ri b_j,
        \qquad
        |x_j|=(a_j^2+b_j^2)^{1/2}.
\end{equation}
Set
\begin{equation}
        \eta(x)=\log\mathfrak p(x)
\end{equation}
and $\Sigma=\{x_j=0:\ \text{ for every } j\in S\}$, the set of possible singularities of $\eta$. By \eqref{asslemma}, all
variables causing these singularities occur with positive power in $x^n$.\\
Let $m$ be a real multi-index in the variables
$(a_1,b_1,\ldots,a_d,b_d)$, and let $m_j$ denote the number of derivatives falling on
the pair $(a_j,b_j)$. Since $\nu_j>r+1$, the functions $|x_j|^{\nu_j}$ are
$C^{r+1}$ across $x_j=0$. Moreover, for $j\in S$ and $1\leq q\leq r+1$,
\begin{equation}\label{der_power_modified}
        |\partial^q |x_j|^{\nu_j}|
        \le
        C_q |x_j|^{\nu_j-q}.
\end{equation}
Since $p_j |x_j|^{\nu_j}\leq \mathfrak p(x)$, this gives, whenever $x_j\neq0$,
\begin{equation}\label{der_power_rho_modified}
        |\partial^q |x_j|^{\nu_j}|
        \le
        C_q \mathfrak p(x)|x_j|^{-q},
        \qquad j\in S.
\end{equation}
At points with $x_j=0$ the derivatives up to order $r+1$ extend continuously because $\nu_j>r+1$; the inverse-power estimate is used only away from the corresponding coordinate hyperplane.
Faà di Bruno's formula applied to $\eta=\log\mathfrak p$ yields, away from $\Sigma$,
\begin{equation}\label{eta_derivative_bound_modified}
        |\partial^m\eta(x)|
        \le
        C_m\prod_{j\in S}|x_j|^{-m_j},
        \qquad
        1\leq |m|\leq r+1.
\end{equation}
For $m=0$, we use the elementary logarithmic bound
\begin{equation}\label{eta_log_bound_modified}
        |\eta(x)|
        \le
        C_0\left(
        1+\sum_{j\in S}|\log |x_j||
        \right).
\end{equation}
Let
\begin{equation}
        u(x)=x^n Q(\eta(x)),
        \qquad
        M=\deg Q.
\end{equation}
By \eqref{eta_log_bound_modified},
\begin{equation}\label{Q_log_bound_modified}
        |Q(\eta(x))|
        \le
        C_Q
        \left(
        1+\sum_{j\in S}|\log |x_j||
        \right)^M.
\end{equation}
Since $n_j>0$ for every $j\in S$, powers dominate logarithms, and therefore
$u(x)\to0$ as $x\to\Sigma$. Thus $u$ extends continuously to $\Sigma$ by setting
$u=0$.\\
We now prove $C^{r,1-\varepsilon}$-regularity in a neighborhood of the origin. Let
$|m|\leq r$. By Leibniz' rule, the chain rule, and
\eqref{eta_derivative_bound_modified}, every term in $\partial^m u$ is bounded by a
finite sum of expressions of the form
\begin{equation}\label{term_bound_modified}
        C_m
        \left(
        \prod_{j\in S}|x_j|^{n_j-m_j}
        \right)
        \left(
        1+\sum_{j\in S}|\log |x_j||
        \right)^M
        R_m(x),
\end{equation}
where $R_m$ is bounded near the origin. Since
\begin{equation}
        |m|\leq r\leq \min_{j\in S}n_j,
\end{equation}
all exponents $n_j-m_j$ are nonnegative. Moreover, by \eqref{regularity_index},
\begin{equation}
        |m|\leq r\leq \sum_{j\in S}n_j-1,
\end{equation}
and hence
\begin{equation}\label{positive_remaining_power_modified}
\sum_{j\in S}(n_j-m_j)\geq 1.
\end{equation}
In particular, at least one positive power of a variable from $S$ remains in every
derivative of order at most $r$. Therefore \eqref{term_bound_modified} implies that each
$\partial^m u$, $|m|\leq r$, extends continuously to $\Sigma$.\\
Let $|m|=r$ and set
\begin{equation}
        \rho(x)=\sum_{j\in S}|x_j|.
\end{equation}
Since \eqref{positive_remaining_power_modified} gives one remaining power, the dominance of powers over logarithms at the origin implies that, for every $\alpha\in(0,1)$,
\begin{equation}\label{pointwise_holder_decay_modified}
        |\partial^m u(x)|
        \le
        C_\alpha \rho(x)^\alpha.
\end{equation}
Differentiating once more away from $\Sigma$, the same estimates give
\begin{equation}\label{gradient_holder_bound_modified}
        |\nabla\partial^m u(x)|
        \le
        C_\alpha \rho(x)^{\alpha-1}.
\end{equation}
If $|x-y|\geq \rho(x)/2$ or $|x-y|\geq \rho(y)/2$, then
\eqref{pointwise_holder_decay_modified} gives
\begin{equation}
        |\partial^m u(x)-\partial^m u(y)|
        \le
        C_\alpha |x-y|^\alpha.
\end{equation}
Otherwise, $|x-y|<\rho(x)/2$ and $|x-y|<\rho(y)/2$. Since $\rho$ is $\sqrt{|S|}$-Lipschitz, after subdividing the segment into a fixed number (depending only on $|S|$) of pieces if necessary, $\rho$ stays comparable to $\rho(x)$ on each piece
and \eqref{gradient_holder_bound_modified} gives
\begin{equation}
        |\partial^m u(x)-\partial^m u(y)|
        \le
        C_\alpha |x-y|\rho(x)^{\alpha-1}
        \le
        C_\alpha |x-y|^\alpha.
\end{equation}
Thus $u\in C^{r,\alpha}$ near the origin for every $\alpha\in(0,1)$. Taking
$\alpha=1-\varepsilon$ gives $u\in C^{r,1-\varepsilon}$.\\
\end{proof}






In the following, we collect some counterexamples as remarks that show that the assumptions of Lemma \ref{lemma_regularity} cannot be relaxed in general. 

\begin{remark}
Functions of the form $x\mapsto x^n Q(\log\mathfrak{p}(x))$ are generally not $C^{N,\alpha}$ for $N\geq |n|$. Indeed, along the real line $x(t)=(t,\dots,t)$, $t>0$, we have that 
\begin{equation}
    \mathfrak{p}(x(t))=t^{\nu_{\min}}(c+o(1)),\qquad
x(t)^n=t^{|n|}, 
\end{equation}
for some constant $c>0$ and thus
\begin{equation}
x^n\log\mathfrak{p}(x(t))=\nu_{\min} t^{|n|}\log t+O(t^{|n|}).
\end{equation}
The function $t^{|n|}\log t$ is $C^{|n|-1}$ but not $C^{|n|}$ at $0$.  The regularity guaranteed by Lemma \ref{lemma_regularity} is thus sharp for non-constant Q, or, more precisely, whenever the logarithmic polynomial has a nonzero logarithmic part.
%We conclude that if $N\ge1$ and $f(0)=0$, then $x^n\log p(x)\in C^{N-1}$but generally not in $C^{N}$ at the origin.
\end{remark}

\begin{remark}
The lower bound on the exponents $\nu_j$ in \eqref{ass_nu_large} cannot be removed in general if
one wants H\"older exponents arbitrarily close to one. Consider the real two-dimensional example
\begin{equation}
        n=(1,1),
        \qquad
        \mathfrak p(x,y)=|x|^\nu+|y|,
        \qquad
        u(x,y)=xy\log\mathfrak p(x,y),
\end{equation}
for $0<\nu<1$. The support condition \eqref{asslemma} is satisfied and the
regularity index predicted by the preceding lemma is $r=1$. For $x>0$, we have
\begin{equation}
        \partial_x u(x,y)
        =
        y\log(x^\nu+|y|)
        +
        \nu y\frac{x^\nu}{x^\nu+|y|}.
\end{equation}
At $x=0$, the continuous extension of this derivative, if it exists, must be
\begin{equation}
        \partial_x u(0,y)
        =
        y\log |y|.
\end{equation}
Fixing $y=y_0>0$, we obtain as $x\to0^+$
\begin{equation}
\begin{split}
        \partial_x u(x,y_0)-\partial_x u(0,y_0)
        &=
        y_0\log\left(1+\frac{x^\nu}{y_0}\right)
        +
        \nu y_0\frac{x^\nu}{x^\nu+y_0}  \\
        &=
        (1+\nu)x^\nu+o(x^\nu).
\end{split}
\end{equation}
Consequently, $\partial_xu$ is not $\alpha$-H\"older continuous in the $x$-direction
for any $\alpha>\nu$. In particular, choosing $\varepsilon>0$ so small that $1-\varepsilon>\nu$, we see that $u\notin C^{1,1-\varepsilon}$.  Thus some lower bound assumption such as \eqref{ass_nu_large} on the exponent $\nu_j$ is necessary for the conclusion $C^{r,1-}$, see also Figure \ref{small_nu_plot}. \\
\begin{figure}
    \centering
\includegraphics[width=0.45\linewidth]{ph.png}
\includegraphics[width=0.45\linewidth]{phx.png}
    \caption{The function $u(x,y) = xy \log(|x|^{\frac{1}{2}}+|y|)$ (blue) compared to its $x$-derivative (orange). The $x$-derivative forms a cusp at $x=0$ which prohibits H\"{o}lder continuity beyond $\alpha = 1/2$.}
    \label{small_nu_plot}
\end{figure}
\end{remark}

\begin{remark}
The support condition in Lemma~\ref{lemma_regularity} is also necessary. Indeed, take
\begin{equation}
        n=(1,0),
        \qquad
        \mathfrak p(x,y)=|y|,
        \qquad
     u(x,y)=x\log|y|.
\end{equation}
Here, the logarithmic power function contains the variable $y$ but the monomial multiplying the logarithmic expression
contains no compensating power of $y$. Along the curve
\begin{equation}
        x=t,
        \qquad
        y=e^{-1/t},
        \qquad
        t>0,
\end{equation}
we have
\begin{equation}
        (x,y)\to(0,0),
        \qquad
        u(t,e^{-1/t})
        =
        t\log(e^{-1/t})
        =
        -1.
\end{equation}
Hence $u$ does not even extend continuously to the origin. This shows that variables
appearing in the logarithmic basis must also appear with positive exponent in the monomial
prefactor.
\end{remark}

\begin{remark}
The regularity in Lemma~\ref{lemma_regularity} cannot, in general, be improved from the
minimal coordinate regularity to the total-degree regularity $|n|-1$. Consider
\begin{equation}
        n=(1,1,1),\quad \mathfrak p(x,y,z)
        =
        x^{100}+y^{10000}+z^{10000},
        \qquad
        u(x,y,z)
        =
        xyz\log\mathfrak p(x,y,z).
\end{equation}
for $x,y,z>0$. The support condition \eqref{asslemma} is satisfied, since every variable appearing in
$\mathfrak p$ also appears in the monomial $xyz$. Moreover, the exponents are
large enough. However,
\begin{equation}
        \min_{j:p_j>0} n_j=1, \text{ while }
        \qquad
        |n|-1=2.
\end{equation}
Thus Lemma~\ref{lemma_regularity} guarantees at most
\begin{equation}
        u\in C^{1,1-\varepsilon},
        \qquad
        \varepsilon\in(0,1),
\end{equation}
whereas a statement with regularity $|n|-1$ would incorrectly predict $C^2$-regularity.\\
Indeed,
\begin{equation}
        \partial_x^2 u(x,y,z)
        =
        yz\left(
        \frac{10100x^{99}}{\mathfrak p(x,y,z)}
        -
        \frac{10000x^{199}}{\mathfrak p(x,y,z)^2}
        \right).
\end{equation}
Along the curve
\begin{equation}
        x=t,
        \qquad
        y=z=t^{1/10},
        \qquad
        t>0,
\end{equation}
we have
\begin{equation}
        \mathfrak p(t,t^{1/10},t^{1/10})
        =
        t^{100}+2t^{1000}
        \sim
        t^{100}.
\end{equation}
Consequently,
\begin{equation}
        \partial_x^2 u(t,t^{1/10},t^{1/10})
        \sim
        C t^{1/5}t^{-2}
        =
        C t^{-9/5},
\end{equation}
which blows up as $t\to0$. Hence $u\notin C^2$. 
This shows that, for mixed monomials, the total degree $|n|$ does not determine the
overall regularity. The limiting quantity is indeed the smallest power with which the
logarithmically active variables occur in the monomial.
\end{remark}


\begin{remark}
    %If we drop assumption \eqref{asslemma}, the Lemma obviously fails, as illustrated by the function \begin{equation}       x \log (x^2+y^2),   \end{equation}which is not continuously differentiable at the origin. Let us stress that the above Lemma also applies for $|n|=1$, since, the function $x\mapsto  x\log|x|$ is $C^{0,\alpha}$ for every $0<\alpha<1$.\\
    From a purely algebraic point of view, the subsequent analysis can be carried out equally with the basis functions
\begin{equation}\label{eta_mod_xj}
    \eta_j = \log|x|_j,\quad 1\leq j\leq d,
\end{equation}
which, in contrast to the logarithms of power functions discussed in Lemma \ref{lemma_regularity}, only guarantee $C^{r-1,1-}$-regularity of functions of the form $x^r_j\log|x_j|$, see Figure \ref{complog}.  Indeed, the basis functions $\log|x_j|$ are unbounded on a whole hyperplane of co-dimension one, where $|x_j|$ is zero. Basis functions given by power functions of the form \eqref{defpower} thus guarantee improved regularity as compared to basis functions of the form \eqref{eta_mod_xj}.



\begin{figure}[ht]
    \centering
\includegraphics[width=0.45\linewidth]{log_derivative_bounded.png}
\includegraphics[width=0.45\linewidth]{log_derivative_unbounded.png}
    \caption{The $x$-derivative of the function $f(x,y) = \log(x^4+y^2)xy$ (blue) compared to the $x$-derivative of the function $g(x,y) = \log(|xy|)xy$ (red). }
    \label{complog}
\end{figure}


\end{remark}



%\todo[inline, color = green]{The above considerations may be extended from power functions to generalized polynomial expression of the form\begin{equation}    \sum_{|m|=1} ^ M p_m(x^\nu_{m})\end{equation}\begin{equation}    \sum_{|m|=1} ^ M p_m x^\nu_{m}\end{equation}}






%\begin{remark}    One might consider adapted norms , then the operator norm estimate \eqref{boundSavinvfull} might imply stronger bounds on as compared to , one might want to consider different topologies, this will be addressed in future works. \end{remark}





%\begin{lemma}Let $x=(x_{i,j})_{1\leq i\leq s, 1\leq j\leq l_i}$ and assume that the matrix $A$ acts on $x$ as $Ax=(\lambda_ix_{i,j}+)$\end{lemma}






\subsection{Approximate Solutions to the Invariance Equation}\label{approx_solutions}

%In the expansion, we either have to use eigen-directions or all directions associated to an eigenvalue, i.e., the whole Jordan block as coordinates to guarantee the invariant action of the linearization. 

We will now construct approximate solutions to the invariance equation \eqref{invariance} by means of logarithmic polynomials. Lemma \ref{shifted} will be used to guarantee that we can solve the cohomology equation order-by-order for specific logarithmic polynomials.
%In order to guarantee regularity of the approximate solution to the invariance equation as a logarithmic polynomial, however, we have to make sure that no mixing of logarithmic terms with monomials of different variables occurs. 


\begin{proposition}\label{approx}
\quad 
Consider an analytic map $f:\mathbb{C}^s\to\mathbb{C}^s$ with $f(0)=0$ and assume that
 \begin{equation}
A=Df(0),
\end{equation}
is semisimple and invertible, i.e.,
\begin{equation}\label{zeronotinspectrum}
0\notin \sigma(A).
\end{equation}
Let $\{\lambda_1,...,\lambda_s\}$ be the set of eigenvalues of $A$ and let $\{e_1,..,e_s\}$ be a basis of eigenvectors of $A$. We assume that the coordinates in $\mathbb{C}^s$ are with respect to the chosen eigenbasis and no generalized eigenvectors are
invoked under the semisimplicity hypothesis.\\
Let $X_1$ be an invariant linear subspace of $A$, i.e., $AX_1\subseteq X_1$, and let $\{e_1,..,e_d\}$ be a basis of $X_1$ consisting of eigenvectors associated to the eigenvalues $\{\lambda_1,...,\lambda_d\}$, counting the eigenvalues up to renumbering, such that
\begin{equation}\label{lambdanotone}
|\lambda_j|< 1,\quad 1\leq j\leq d.
\end{equation} 
For any $N\in\mathbb{N}$, we can find a logarithmic polynomial of degree not larger than $N$,
\begin{equation}\label{hN}
h^{< N}(x)=\sum_{1\leq |n|< N} h_n(\eta)x^n,
%h^{< N}=\sum_{1\leq |n|< N} h_n(\eta_1,...,\eta_M)x^n,
\end{equation}
where $\eta = (\eta_1,...,\eta_M)$ is a vector of specific logarithm power functions of the form \eqref{def_eta} such that the support restrictions of Lemma \ref{lemma_regularity} are satisfied. The polynomials $h_n:\mathbb{C}^M\to\mathbb{C}^s$ are
normalized to
\begin{equation}\label{pO1}
h_{1@l}=v_l, \quad 1\leq l\leq d,
\end{equation}
and satisfy 
\begin{equation}\label{bounddegree}
    \deg h_n \leq |n|-1,
\end{equation}
such that the invariance equation \eqref{invariance} can be solved up to order $N$:
\begin{equation}
f(h^{< N}(x))=h^{< N}(Ax)+R_N(x),
\end{equation}
where the remainder is given as a logarithmic polynomial of the form
\begin{equation}
R_N(x)=\sum_{|n|\geq N}r_n(\eta)x^n, %_N(x)=\sum_{|n|\geq N}r_n(\eta_1,...,\eta_M)x^n, 
\end{equation}
for polynomials $r_n:\mathbb{C}^M\to\mathbb{C}^s$, satisfying the support conditions of Lemma \ref{lemma_regularity} as well. 
\end{proposition}


\begin{proof}
%To ease the traceability of the proof, we treat the \rev{semisimple} and the non-\rev{semisimple} case separately. This also provides a more explicit description of the \rev{semisimple} case, which often appears in applications. In this case, we have that $s=d$ and that $\eta_j=x_j$ for $1\leq j\leq s$. 

We write $\lambda=(\lambda_1,...,\lambda_s)$ for the vector of eigenvalues. Since the coordinates are with respect to eigenvectors, we have that
\begin{equation}\label{Ax}
A|_{X_1}x=(\lambda_1x_1,..,\lambda_dx_d).
\end{equation}
The map $f$ being analytic, it admits a Taylor series expansion up to any order and we can write 
\begin{equation}
f(x)=Ax+\sum_{|n|= 2}^\infty f_n x^n,
\end{equation}
and we can expand $f(h(x))$ in powers of $x$ with coefficients depending on $\eta_j$, $1\leq j\leq M$: 
\begin{equation}\label{expand_f}
\begin{split}
f(h(x))&=\sum_{1\leq |n|\leq N} Ah_n(\eta)x^n+\sum_{|n|= 2}^r f_n \left(\sum_{1\leq |m|\leq N} h_m(\eta)x^m\right)^n+\mathcal{O}(|x|^{r+1}).
\end{split}
\end{equation}
Comparing powers of $x$ in $f(h(x))-h(Ax)$ and using \eqref{Ax}, we find that the invariance equation at $x^n$ becomes
\begin{equation}\label{coholn}
Ah_n(\eta(x))-\lambda^nh_n(\eta(Ax))=Q_n,
\end{equation}
where $Q_n$ is a logarithmic polynomial depending on the $h_m$'s for $|m|<|n|$. The support condition \eqref{asslemma} is preserved under the operations performed in \eqref{expand_f}: multiplication adds the monomial multi-indices and the logarithmic multiplicities, addition preserves the condition term by term, and the shift $\eta\mapsto\eta+v$ can only lower logarithmic powers. Thus, inductively, each logarithmic factor introduced at a resonant multi-index remains accompanied by the corresponding monomial support, and $Q_n$ satisfies \eqref{asslemma} as well. 
In particular, for $n=1@j$, $1\leq j\leq d$, we have that
\begin{equation}
Ah_{1@j}-\lambda_jh_{1@j}=0,
\end{equation}
which holds by assumption \eqref{pO1}.\\
We proceed inductively in solving equation \eqref{coholn}. To this end, we expand the polynomial $h_n$ in the eigenbasis $\{e_1,..e_s\}$ and write
\begin{equation}
    h_{n,j}(x)=\langle h_n(x),e_j\rangle,\quad 1\leq j\leq s. 
\end{equation}
Similarly, we write
\begin{equation}
    Q_{n,j}(x)=\langle Q_n(x),e_j\rangle,\quad 1\leq j\leq s.
\end{equation}
Since the components of $\eta$ are power functions of the form \eqref{def_eta}, they satisfy
\begin{equation}
    \eta_l(Ax) = \eta_l(x) + v_l,\quad 1\leq l\leq M,
\end{equation}
%\begin{equation}\Lambda=(\log|\lambda_1|,...,\log|\lambda_d|),\quad \eta=(\eta_1,...,\eta_M),\end{equation}
and equation \eqref{coholn} simply becomes 
\begin{equation}\label{eqpnj}
\lambda_jh_{n,j}(\eta)-\lambda^nh_{n,j}(\eta+v)=Q_{n,j}(\eta),\quad 1\leq j\leq s,
\end{equation}
for some shift vector $v=(v_1,...,v_M)$. 
Up to the first resonant index, the right-hand side of \eqref{eqpnj} are constants, i.e., logarithmic polynomials of degree zero and we can subsequently solve the shifted polynomial equation using Lemma \ref{shifted} for logarithmic polynomials $h_n$ of degree zero as well.\\
Let us take a closer look at what happens at one of the first resonant indices $n \in I_{\rm min}(\Lambda)$. In fact, there might be several resonant indices of minimal order and we may pick one of them for illustration:
\begin{equation}\label{eqfirstresonance}
    h_{n,j}(\eta) - h_{n,j}(\eta+v) = -(\lambda_j)^{-1} Q_{n,j},
\end{equation}
where we have used that $\lambda^n=\lambda_j$ for some j and $Q_{n,j}$ is still a constant. By Lemma \ref{shifted}, we know that any first-order polynomial $h_{n,j}(\eta) = h_{n,j,0}+h_{n,j,1}\eta$ with $ v\cdot h_{n,j,1} = (\lambda_j)^{-1} Q_{n,j} $ defines a solution to \eqref{eqfirstresonance}.
%Only those polynomials, however, which satisfy Condition \eqref{condition} guarantee that $p_{n_{\rm min}}x^{n_{\rm min}}$ is $C^\alpha$ for any $\alpha\in (0,1)$. The non-zero entries of $\pi_1$ thus correspond to those entries of $n_{\rm min}$ which are non-zero. Condition \eqref{condition} is preserved under taking powers and thus any higher-order terms satisfy the condition as well.
If the index $n$ is non-resonant, we can apply Lemma \ref{shifted} and find a polynomial solution to equation \eqref{eqpnj} of the same degree as $\deg Q_{n,j}$. If, on the other hand, the index $n$ is resonant, we can apply Lemma \ref{shifted} and find a polynomial solution to equation \eqref{eqpnj} of degree $\deg Q_{n,j}+1$.\\
%in the same variables as $Q_{n,j}+1$. This proves the claim. Condition \eqref{condition} is thus propagated to any order of the approximate solution to the invariance condition.\\
Finally, to prove the bound \eqref{bounddegree}, we proceed by induction over the order of the multi-index $n$. At order one, the bound is trivially satisfied by the normalization \eqref{pO1}. Assume now that the claim has been proven for any multi-index up to order $|n|$.\\
At order $k=|n|+1$, every nonlinear contribution to the right-hand side of \eqref{eqpnj} is a product of at least two previously constructed terms of orders $k_1,\ldots,k_j$ with $j\geq2$ and $k_1+\cdots+k_j=k$. By the induction hypothesis its logarithmic degree is therefore at most
\begin{equation}
(k_1-1)+\cdots+(k_j-1)=k-j\leq k-2=|n|-1.
\end{equation}
If the index is non-resonant, Lemma \ref{shifted} gives a solution of the same degree, hence degree at most $k-2$. If the index is resonant, the degree is raised by at most one, hence is at most $k-1$, which is exactly \eqref{bounddegree}.
\end{proof}


\begin{corollary}
Let $h^{<N}$ be a logarithmic polynomial of the form \eqref{hN}, solving the invariance equation up to order $N$ and let $r_*\geq 1$ be defined as in \eqref{def_r_star}. Then the approximate solution $h^{<N}$ is $C^{r_*,1-}$.
\end{corollary}

\begin{proof}
   This is an immediate consequence from Proposition \ref{approx} and Lemma \ref{lemma_regularity}. The power functions defining $\eta_l$ have to be chosen such that the support of the weights matches the non-zero entries of the resonant indices, see assumption \eqref{asslemma}, while the shifts $v_l$ have to be chosen in accordance with the largeness assumption \eqref{ass_nu_large} on the exponents $\nu_j$. This can always be achieved choosing $\gamma$ in \eqref{nu_form} sufficiently small. 
\end{proof}

%\begin{remark}    The index $\alpha$ in the norm just tells us that we can control the logarithmic polynomial appearing at the lowest order in $x^n$ of the remainder. \end{remark}

%The logarithmic polynomials appearing in the approximate solution to the invariance equation have more structure.

%\begin{remark}    Derivatives are not necessarily zero. The regularity of the manifold is determined by the appearance of the first resonance. \end{remark}




%\begin{remark}Let us remark that, even in the presence of resonances, logarithmic contributions only enter the approximation \eqref{hN}, i.e., $p_n\neq 1$, if the  corresponding resonant Taylor coefficient of the  map $f$ is non-zero.\end{remark}

%\begin{proposition}assume that \textcolor{red}{$f:\mathbb{R}^{2n}\to\mathbb{R}^{2n}$ is $f|_{\mathbb{R}^n}$ is r-times continuously differentiable with}  $r\geq 2$ \todo{r=1?}(i.e. admits admits a Taylor expansion up to order $r$)\end{proposition}










\section{Proof of the Existence of Resonant Invariant Manifolds}\label{existence}

In this section, we pass from the finite-order logarithmic approximation constructed above to an exact invariant manifold. The argument has two logically separate parts. First, after all resonant orders have been incorporated into the approximate solution, the remaining tail starts at an order $N$ beyond every resonance on which the linear conjugacy operator is invertible with a uniform bound. Second, the nonlinear remainder is small on a sufficiently small neighborhood, so the corrected invariance equation is a contraction. This is the analytic step that turns the algebraic approximate solution into a genuine one. \\
Traditionally, the convergence of the Taylor expansion of analytic invariant manifolds about an equilibrium is proved using the method of majorants \cite{moser1956analytic,moser2016stable}, a technique that traces back to the pioneering works of Hadamard \cite{hadamard1901} and Perron \cite{perron1929stabilitat}. A key ingredient of this approach is the simple yet crucial observation that the coefficients appearing in the Faa di Bruno formula are nonnegative \cite{comtet2012advanced,siegel2012lectures}. More recently, modern treatments reformulate the invariance equation as a fixed-point problem and establish existence results in suitably chosen function spaces \cite{figueras2017framework,hirsch1970invariant}.\\


\subsection{Invertibility of the Linear Conjugacy Operator}

For an $s\times s$ semisimple matrix $A$ and an $A$-invariant $d$-dimensional subspace $X_1\subseteq \mathbb{C}^s$, define the operator $\mathcal{S}_A$, acting on functions $H:B^1_{d}\to\mathbb{C}^s$, $B^1_{d}\subset X_1$, as
	\begin{equation}\label{defS}
	\mathcal{S}_AH(x):=AH(x)-H(A|_{X_1}x).
	\end{equation}
The operator $\mathcal{S}_A$ is the linearization, with respect to the correction $H$, of the invariance-linearization equation around the linear dynamics and defines the vector-valued analogue of \eqref{defSav}. Its properties on the Banach space $\Gamma_N^d$ are summarized in the following lemma.



\begin{proposition}\label{Lemma_SM}
Let $A\in\mathbb C^{s\times s}$ be semisimple and invertible and let $\lambda_1,\ldots,\lambda_s$ denote the eigenvalues of $A$, counted with
multiplicity. Let
$X_1\subset\mathbb C^s$ be a $d$-dimensional $A$-invariant subspace associated with
eigenvalues $\lambda_1,\ldots,\lambda_d$, and assume
\begin{equation}
        |\lambda_i|<1,
        \qquad
        1\leq i\leq d.
\end{equation}
Suppose $N$ is chosen so large that
\begin{equation}\label{nonres_tail}
        \lambda^n\ne \lambda_j,
        \qquad
        |n|\geq N,
        \qquad
        1\leq j\leq s.
\end{equation}
Then, for $\beta>0$ sufficiently small, the operator
\begin{equation}
        \mathcal S_A:\Gamma_N^d\to\Gamma_N^d,
        \qquad
        \mathcal S_A H(x)=AH(x)-H(Ax),
\end{equation}
is bounded and invertible. Moreover, there exists $q_*\in(0,1)$ such that
\begin{equation}\label{SA_inverse_bound}
        \|\mathcal S_A^{-1}\|_{\rm op}
        \le
        \frac{1}{1-q_*}
        \frac{1}{\min_{1\leq j\leq s}|\lambda_j|}.
\end{equation}
\end{proposition}

\begin{proof}
Choose coordinates on $X_1$ along eigenvectors of $A|_{X_1}$, and coordinates in
$\mathbb C^s$ along eigenvectors of $A$. Thus
\begin{equation}
        Ax=(\lambda_1x_1,\ldots,\lambda_dx_d),
        \qquad
        AH_n=(\lambda_1H_{n,1},\ldots,\lambda_sH_{n,s}).
\end{equation}
Let
\begin{equation}
        H(x)=\sum_{|n|\geq N}H_n(\eta)x^n,
        \qquad
        \widetilde H(x)=\sum_{|n|\geq N}\widetilde H_n(\eta)x^n.
\end{equation}
We consider the equation
\begin{equation}
        \mathcal S_AH=\widetilde H
\end{equation}
and assume again that the basis functions have been chosen such that $\eta(Ax)=\eta(x)+v$. The coefficient equation then becomes
\begin{equation}\label{coeff_SA}
        \lambda_j H_{n,j}(\eta)
        -
        \lambda^n H_{n,j}(\eta+v)
        =
        \widetilde H_{n,j}(\eta),
        \qquad
        |n|\geq N,
        \qquad
        1\leq j\leq s,
\end{equation}
or, equivalently,
\begin{equation}\label{coeff_shift}
        H_{n,j}(\eta)
        -
        a_{n,j}H_{n,j}(\eta+v)
        =
        \lambda_j^{-1}\widetilde H_{n,j}(\eta),
        \qquad
        a_{n,j}=\frac{\lambda^n}{\lambda_j}.
\end{equation}
We now choose $N$ and $\beta$ to guarantee invertibility of $\mathcal{S}_A$. Since
\begin{equation}
        \max_{1\leq i\leq d}|\lambda_i|<1,
\end{equation}
there exist $N$ sufficiently large and $\beta>0$ sufficiently small such that
\begin{equation}\label{qstar_condition}
        q_*
        :=
        \sup_{\substack{|n|\geq N\\1\leq j\leq s}}
        \left|
        \frac{\lambda^n}{\lambda_j}
        \right|
        e^{\beta |v||n|}
        <1.
\end{equation}
Indeed,
\begin{equation}
        \left|
        \frac{\lambda^n}{\lambda_j}
        \right|
        e^{\beta |v||n|}
        \le
        \frac{
        \left(\max_{1\leq i\leq d}|\lambda_i|\right)^{|n|}
        e^{\beta |v||n|}
        }{
        \min_{1\leq j\leq s}|\lambda_j|
        },
\end{equation}
and the right-hand side tends to zero as $|n|\to\infty$ if $\beta>0$ is chosen so small that
\begin{equation}
        e^{\beta |v|}\max_{1\leq i\leq d}|\lambda_i|<1.
\end{equation}
By \eqref{nonres_tail}, $a_{n,j}\ne1$ for all $|n|\geq N$ and $1\leq j\leq s$.
The shifted-polynomial estimate applied to \eqref{coeff_shift}, with shift vector $v$, gives
\begin{equation}\label{coeff_inverse_bound}
        \|H_{n,j}\|_{\beta |n|}
        \le
        \frac{1}{1-|a_{n,j}|e^{\beta |v||n|}}
        \frac{1}{|\lambda_j|}
        \|\widetilde H_{n,j}\|_{\beta |n|}.
\end{equation}
Using \eqref{qstar_condition}, we obtain
\begin{equation}
        \|H_{n,j}\|_{\beta |n|}
        \le
        \frac{1}{1-q_*}
        \frac{1}{\min_{1\leq j\leq s}|\lambda_j|}
        \|\widetilde H_{n,j}\|_{\beta |n|}.
\end{equation}
Summing over $|n|\geq N$ and over $1\leq j\leq s$ yields
\begin{equation}
        \|H\|_{\Gamma_N^d,\beta,\sigma}
        \le
        \frac{1}{1-q_*}
        \frac{1}{\min_{1\leq j\leq s}|\lambda_j|}
        \|\widetilde H\|_{\Gamma_N^d,\beta,\sigma}.
\end{equation}
Therefore $\mathcal S_A^{-1}$ exists and satisfies \eqref{SA_inverse_bound}. Boundedness
of $\mathcal S_A$ follows directly from the definition of the norm, the finite-dimensional
boundedness of $A$, and the shift estimate for $\eta\mapsto\eta+v$. This proves the claim.
\end{proof}



\begin{remark}
    The invertibility of the operator $\mathcal{S}_A$, where $A$ is defined in \eqref{defA}, will be crucial in the reformulation of the invariance equation as a fixed-point problem in the next section. To guarantee invertibility, we have to assume that the approximation constructed in Proposition \ref{approx} is such that it includes all resonant indices.
    %This is necessary to guarantee the third assumption in Lemma \ref{LemmaSM}. 
    In particular, once an approximation is fixed, the remainder constructed by the fixed point argument will be unique. We further note that all the bounds used in the proof of Proposition \ref{Lemma_SM} can be reformulated in terms of the matrix norms of $A$ and $A|_{X_1}$ instead of their eigenvalues. Indeed, thanks to the Gelfand Theorem, we may always choose an adapted norm which is arbitrarily close to the spectral radius of both matrices. 
    %This rules out the possibility that resonances with larger index could violate the uniqueness of the constructed invariant manifold. 
\end{remark}







\subsection{Reformulation of the Invariance Equation as a Fixed-Point Problem}

We will assume that $x\mapsto h(x)$ is defined on a small ball $B_{\delta,d}^\mathbb{C}$, which, after re-scaling $x\mapsto \delta x$ and setting $h^\delta(x)=h(\delta x)$, amounts to the re-scaled invariance equation:
\begin{equation}\label{coholdelta}
Ah^\delta(x)+N_f(h^\delta(x))=h^\delta(Ax).
\end{equation}
Dividing equation \eqref{coholdelta} by $\delta$ and setting $h(x)=h^\delta(x)/\delta$ by abuse of notation, the invariance equation becomes
\begin{equation}\label{invariance_expand_f}
Ah(x)+\delta N_f^\delta (h(x))=h(Ax),
\end{equation}
where
\begin{equation}\label{N_f_expand}
N_f^\delta(x)=\sum_{|n|=2}^\infty \delta^{|n|-2} f_nx^n,
\end{equation}
is the re-scaled nonlinear part of $f$. In particular, we can assume, by choosing $\delta$ small, that the nonlinear part of the invariance equation is as small as we like in an appropriately chosen norm, while the parametrization $h$ is defined on the unit ball $B_{1,d}^\mathbb{C}$.\\
Assume now that we have found an approximate solution $h^{< N}$ of order $N$, according to Proposition \ref{approx} and Proposition \ref{Lemma_SM}. In the presence of resonances, where there is no longer a unique approximate solution to the invariance equation, any member of the family of approximate solutions will do. We write
\begin{equation}
h(x)=h^{< N}(x)+H(x),
\end{equation}
for an unknown function $H$. Since the approximation is a logarithmic polynomial of degree $N-1$, the remainder $H$ is assumed to be a logarithmic polynomial with leading polynomial contribution of order $N$.
%Since $|x|^{-N-1}H(x)$ contains, in general, logarithmic terms, we can only guarantee that\todo[inline, color = red]{Include leading order behavior of $H$}
We can split the invariance equation \eqref{invariance} as
\begin{equation}
f(h^{< N}(x)+H(x))=h^{< N}(Ax)+H(Ax),
\end{equation}
which, after expanding the map $f$ according to \eqref{invariance_expand_f},
\begin{equation}\label{invariance_rewritten}
    A h^{<N}+ AH + \delta N_f^\delta\circ(h^{< N}+H) = h^{< N}\circ A+H\circ A.
\end{equation}
Regrouping and using the definition of the operator $\mathcal{S}_A$ defined in \eqref{defS}, equation \eqref{invariance_rewritten} reads 
\begin{equation}\label{coholH}
\mathcal{S}_AH = h^{< N}\circ A-Ah^{< N} -\delta N_f^\delta\circ(h^{< N}+H). 
\end{equation}
By Proposition \ref{Lemma_SM}, we can invert the operator $\mathcal{S}_A$ in the space $\Gamma_{N}^d$ and rewrite equation \eqref{coholH} as the fixed point problem
\begin{equation}
H=\mathcal{T}_\delta(H,h^{< N}),
\end{equation}
for the functional
\begin{equation}\label{defT}
\mathcal{T}_\delta(H,h^{< N}) =\mathcal{S}_A^{-1}\Big(h^{< N}\circ A-Ah^{< N} -\delta N_f^\delta(h^{< N}+H)\Big).
\end{equation}
We split the right-hand side into a fixed residual and a nonlinear tail:
\begin{equation}
        \tau_\delta
        =
        h^{<N}\circ A-Ah^{<N}
        -
        \delta N_f^\delta(h^{<N}),
\end{equation}
and
\begin{equation}
        \mathcal N_\delta(H)
        =
        -\delta\left[
        N_f^\delta(h^{<N}+H)-N_f^\delta(h^{<N})
        \right],
\end{equation}
such that the operator \eqref{defT} reads
\begin{equation}
    \mathcal{T}_\delta(H,h^{<N}) =  \mathcal{S}_A^{-1}\left(\tau_\delta+\mathcal N_\delta(H)\right)
\end{equation}




\begin{remark}
    Since the approximate solution $h^{< N}$ is not unique in the presence of resonances, the functional \eqref{defT} will vary for different approximate solutions. Once an approximation is fixed, however, the fixed-point of $\mathcal{T}_\delta$ will be unique. 
\end{remark}

Before we state the main result of this section about fixed-points of the functional \eqref{defT}, corresponding to invariant manifolds conjugated to the linear part, we collect some functional analytic properties of $\mathcal{T}_\delta$. The following lemma is the standard regularity for affine tail composition, see also \cite{hille1996functional} for further details. 

%the \emph{Nemytskii operator} and is included for completeness, see also \cite{hille1996functional} for further details.

\begin{lemma}[Affine tail composition]\label{lemma_affine_tail}
Let $h^{<N}$ be a fixed logarithmic polynomial with
\begin{equation}
        h^{<N}(x)=Lx+O(|x|^2),
\end{equation}
and let $H\in\Gamma_N^d$. Let $N_f^\delta$ be analytic near the origin and satisfy
\begin{equation}
        N_f^\delta(z)=O(|z|^2).
\end{equation}
Then
\begin{equation}\label{N_delta}
        \mathcal N_\delta(H)
        =
        -\delta\left[
        N_f^\delta(h^{<N}+H)-N_f^\delta(h^{<N})
        \right]
\end{equation}
belongs to $\Gamma_N^d$. Moreover, for every $R>0$ sufficiently small, there exists
$C_R>0$ such that, for all $H_1,H_2\in\Gamma_N^d$ with
\begin{equation}
        \|H_i\|_{\Gamma_N^d,\beta,\sigma}\leq R,
        \qquad i=1,2,
\end{equation}
one has
\begin{equation}\label{bound_difference_N_delta}
        \|\mathcal N_\delta(H_1)-\mathcal N_\delta(H_2)\|_{\Gamma_N^d,\beta,\sigma}
        \le
        \delta C_R
        \|H_1-H_2\|_{\Gamma_N^d,\beta,\sigma}.
\end{equation}
\end{lemma}

\begin{proof}
Since $N_f^\delta$ is analytic and vanishes to second order, we may write
\begin{equation}
        N_f^\delta(z)
        =
        \sum_{k\ge2} B_k^\delta(z,\ldots,z),
\end{equation}
where $B_k^\delta$ are symmetric $k$-linear maps. Therefore
\begin{equation}
        N_f^\delta(h^{<N}+H)-N_f^\delta(h^{<N})
        =
        \sum_{k\ge2}
        \sum_{\ell=1}^k
        \binom{k}{\ell}
        B_k^\delta(
        \underbrace{H,\ldots,H}_{\ell},
        \underbrace{h^{<N},\ldots,h^{<N}}_{k-\ell}).
\end{equation}
Every term on the right-hand side contains at least one factor $H$. Since
$H\in\Gamma_N^d$ starts at order $N$, and multiplication by the fixed logarithmic
polynomial $h^{<N}$ preserves the tail order $N$, each term belongs to $\Gamma_N^d$.\\
The Banach-algebra property of $\Gamma_N^d$, together with boundedness of multiplication by
the fixed polynomial $h^{<N}$, gives convergence of the above series for
$\|H\|_{\Gamma_N^d,\beta,\sigma}\leq R$, after choosing $R$ and the domain sufficiently small.
Similarly, subtracting the two expansions for $H_1$ and $H_2$, every difference factors
one copy of $H_1-H_2$. Hence
\begin{equation}
        \|\mathcal N_\delta(H_1)-\mathcal N_\delta(H_2)\|_{\Gamma_N^d,\beta,\sigma}
        \le
        \delta C_R
        \|H_1-H_2\|_{\Gamma_N^d,\beta,\sigma}.
\end{equation}
This proves the claim.
\end{proof}








\begin{proposition}\label{propfixedpoint}
Let the assumptions of Theorem \ref{mainthm} and Proposition \ref{Lemma_SM} with
$\beta$ sufficiently small be met. Assume further that the invariance equation has been
scaled according to \eqref{coholdelta} for $\delta>0$ sufficiently small. Then, for any
approximate solution $h^{<N}$, with $N\in\mathbb N$ chosen larger than the order of every resonant index, which is possible by Lemma~\ref{lemma_finite_resonances}) and sufficiently large for Proposition~\ref{Lemma_SM}, the functional $\mathcal T_\delta(\cdot,h^{<N})$ is a contraction from the
closed unit ball $B^1(\Gamma_N^d)$ into itself and thus has a unique fixed point.
\end{proposition}

\begin{proof}
To avoid cluttering the notation, we write
\begin{equation}
        \|\cdot\|=\|\cdot\|_{\Gamma_N^d,\beta,\sigma}.
\end{equation}
By Lemma \ref{lemma_affine_tail}, $\mathcal N_\delta(H)\in\Gamma_N^d$ for
$H\in\Gamma_N^d$, and $\tau_\delta\in\Gamma_N^d$ because $h^{<N}$ solves the scaled
invariance equation up to order $N-1$. Thus $\mathcal T_\delta$ is well-defined as a map
from $\Gamma_N^d$ to itself.

We first prove that $\mathcal T_\delta$ is a contraction on $B^1(\Gamma_N^d)$. Let
$H,\widetilde H\in B^1(\Gamma_N^d)$. By Lemma \ref{lemma_affine_tail}, applied with
$R=1$, there exists a constant $C_1>0$ such that
\begin{equation}\label{contraction_revised}
\begin{split}
        \|\mathcal T_\delta(H,h^{<N})-\mathcal T_\delta(\widetilde H,h^{<N})\|
        &=
        \left\|
        \mathcal S_A^{-1}
        \left(
        \mathcal N_\delta(H)-\mathcal N_\delta(\widetilde H)
        \right)
        \right\| \\
        &\le
        \|\mathcal S_A^{-1}\|_{\rm op}
        \|\mathcal N_\delta(H)-\mathcal N_\delta(\widetilde H)\| \\
        &\le
        \delta
        \|\mathcal S_A^{-1}\|_{\rm op}
        C_1
        \|H-\widetilde H\|.
\end{split}
\end{equation}
Choosing $\delta>0$ sufficiently small so that
\begin{equation}
        \delta
        \|\mathcal S_A^{-1}\|_{\rm op}
        C_1
        <1,
\end{equation}
we obtain that $\mathcal T_\delta(\cdot,h^{<N})$ is a contraction on
$B^1(\Gamma_N^d)$.\\
It remains to show that $\mathcal T_\delta$ maps the closed unit ball into itself. For
$H\in B^1(\Gamma_N^d)$, we have
\begin{equation}\label{ball_revised}
        \mathcal T_\delta(H,h^{<N})
        =
        \mathcal S_A^{-1}\tau_\delta
        +
        \mathcal S_A^{-1}\mathcal N_\delta(H).
\end{equation}
The approximate solution $h^{<N}$ is constructed for the scaled
invariance equation through order $N-1$. Therefore the residual $\tau_\delta$ starts at order $N$, and hence
\begin{equation}
    \tau_\delta\in\Gamma_N^d.
\end{equation}
Moreover, by \eqref{N_f_expand}, the homogeneous term of order
$k$ in the scaled nonlinear part $\delta N_f^\delta$ carries the
factor $\delta^{k-1}$. For completeness, we make the dependence on the scaling parameter $\delta$ explicit. Writing again
\begin{equation}
N_f^\delta(z)
=
\sum_{k\geq 2} B_k^\delta(z,\ldots,z),
\qquad
B_k^\delta=\delta^{k-2}B_k,
\end{equation}
we have, for $0<\delta\leq 1$, the uniform bounds
\begin{equation}
\|B_k^\delta\|\leq \|B_k\|.
\end{equation}
Since $h^{<N}$ is a fixed logarithmic polynomial, multiplication by
$h^{<N}$ defines a bounded operator on $\Gamma_N^d$. Hence, after
restricting the domain if necessary, analyticity of $N_f$ implies that,
for $\|H_i\|_{\Gamma_N^d}\leq R$,
\begin{equation}
\bigl\|
N_f^\delta(h^{<N}+H_1)
-
N_f^\delta(h^{<N}+H_2)
\bigr\|_{\Gamma_N^d}
\leq
C_R\|H_1-H_2\|_{\Gamma_N^d},
\end{equation}
where $C_R$ can be chosen independently of sufficiently small $\delta$.
Multiplication by the prefactor $\delta$ in \eqref{N_delta} therefore gives
\eqref{bound_difference_N_delta}. Moreover, the order-by-order construction of $h^{<N}$ shows inductively that its homogeneous coefficient of order $k$ is of
the form
\begin{equation}
h_k^\delta=\delta^{k-1}\widehat h_k,
\end{equation}
with $\widehat h_k$ independent of $\delta$. Indeed, an order-$k$
contribution arising from an $m$-linear term of $N_f$, evaluated on
homogeneous terms of orders $j_1,\ldots,j_m$ with
$j_1+\cdots+j_m=k$, carries the factor
\begin{equation}
\delta^{m-1}\prod_{i=1}^m\delta^{j_i-1}
=
\delta^{k-1},
\end{equation}
and inversion of the corresponding cohomology operator does not alter
this factor. Since all terms of orders smaller than $N$ cancel by
construction, the residual $\tau_\delta$ therefore starts at order $N$,
and its coefficients of order $k\geq N$ carry a factor $\delta^{k-1}$. The same analytic majorant estimate as above then yields, uniformly for sufficiently small $\delta$,
\begin{equation}\label{bound_residual_revised}
\|S_A^{-1}\tau_\delta\|_{\Gamma_N^d}
\leq
C_\tau\delta^{N-1},
\end{equation}
for some constant $C_\tau>0$. For the nonlinear tail, Lemma \ref{lemma_affine_tail} gives, again with $R=1$,
\begin{equation}
        \|\mathcal N_\delta(H)-\mathcal N_\delta(0)\|
        \le
        \delta C_1\|H\|.
\end{equation}
Since $\mathcal N_\delta(0)=0$, this yields
\begin{equation}\label{tail_bound_revised}
        \|\mathcal S_A^{-1}\mathcal N_\delta(H)\|
        \le
        \delta
        \|\mathcal S_A^{-1}\|_{\rm op}
        C_1
        \|H\|.
\end{equation}
Combining \eqref{ball_revised}, \eqref{bound_residual_revised}, and
\eqref{tail_bound_revised}, we obtain
\begin{equation}
    \|T_\delta(H,h^{<N})\|
    \leq
    C_\tau\delta^{N-1}
    +
    \delta\|S_A^{-1}\|_{\rm op}C_1\|H\|.
\end{equation}
For $H\in B^1(\Gamma_N^d)$, this gives
\begin{equation}
        \|\mathcal T_\delta(H,h^{<N})\|
        \le
        C_{\tau}\delta^{N-1}
        +
        \delta
        \|\mathcal S_A^{-1}\|_{\rm op}
        C_1 .
\end{equation}
By decreasing $\delta>0$ further, if necessary, we can ensure that
\begin{equation}
        C_{\tau}\delta^{N-1}
        +
        \delta
        \|\mathcal S_A^{-1}\|_{\rm op}
        C_1
        \leq 1.
\end{equation}
Thus $\mathcal T_\delta(\cdot,h^{<N})$ maps the closed unit ball
$B^1(\Gamma_N^d)$ into itself.\\
The Banach fixed point theorem now gives a unique fixed point
$H\in B^1(\Gamma_N^d)$. Consequently,
\begin{equation}
        h=h^{<N}+H
\end{equation}
solves the scaled invariance equation. This proves the claim.
\end{proof}



%Since $\Gamma^d_N$ is closed under composition with analytic functions from the left and since $\tilde{H}$ belongs to $\Gamma_{N}^d$, we find that\begin{equation}    DN_f^\delta(h^{< N}+H)\tilde{H}\in \Gamma_{N}^d,\end{equation}as well. Therefore, we have that the increment $\mathcal{T}_\delta(H+\tilde{H},h^{< N})-\mathcal{T}_\delta(H,h^{< N})$ belongs to $\Gamma_{N}^d$ as well. Taking the norm \eqref{defnorm} of the increment formula gives\begin{equation}\label{contraction}\begin{split}\|\mathcal{T}_\delta(H+\tilde{H},h^{< N})-\mathcal{T}_\delta(H,h^{< N})\|_{\Gamma^d_N,\beta} & = \delta\|\mathcal{S}_A^{-1}[N_f^\delta(h^{< N}+H)-N_f^\delta(h^{<N}+\tilde{H})]\|_{\Gamma_N^d,\beta,\sigma}\end{split}\end{equation}




%For $H, \tilde{H}\in \Gamma_{N}^d$, we have the finite increments formula\begin{equation}\label{increment}\begin{split}\mathcal{T}_\delta(H+\tilde{H},h^{< N})-\mathcal{T}_\delta(H,h^{< N})&=\int_0^1 \frac{d}{ds} \mathcal{T}_\delta(H+s\tilde{H},h^{< N})\, ds \\&=\int_0^1\frac{d}{ds}\mathcal{S}^{-1}_A\Big(h^{< N}\circ A-Ah^{< N} -\delta N_f^\delta(h^{< N}+H+s\tilde{H})\Big)\, ds\\&=-\delta\int_0^1\mathcal{S}^{-1}_ADN_f^\delta(h^{< N}+H+s\tilde{H})\tilde{H}\, ds,\end{split}\end{equation}where $DN_f^\delta$ denotes the Jacobian of the nonlinear part $N_f^\delta$.



%, we first note that all derivatives up to order $N$ of the expression $\mathcal{T}_\delta(H,h^{< N})-\mathcal{T}_\delta(0,h^{< N})$ are zero (see the above calculations on the contraction property) and, by

%\begin{lemma} \label{ascoli-arzela}	Assume that $\{u_n\}_{n \in \mathbb{N}} \subset B^{\Gamma_{s,l}}_1$ 	(respectively $\{u_n\}_{n \in \mathbb{N}} \subset B^{\Upsilon_{s,l}}_1$) 	and that for all $x \ne 0 $, $u_n(x) $ converges to $u(x)$ in the $C^0$-norm.	Then $u(x) \in B^{\Gamma_{s,l-1}}_1$ (respectively  $B^{\Gamma_{s,l-1}}_1$).\\	Furthermore,  for each $x \ne 0 $ and any $k < s$ the derivative $D^{k} u_n(x)$ converges to $D^k u(x)$	uniformly on compact sets of $B_{1,d} \setminus \{0\}$. \end{lemma} \begin{proof} 	This is a version of Ascoli-Arzela theorem, see e.g. \cite{}.\\	Given $x \in \mathbb{C}^d$, we consider a ball $B$ centered in $x$ and radius $1/10 |x|$. 	By the Ascoli--Arzela theorem, for any subsequence of $u_n$, we can extract a $C^{l-1}(B)$ 	convergent subsequence. By the assumption that $u_n$ converges pointwise, we obtain 	that the limit should be $u$	\footnote{Indeed, it is easy to see that if any subsequence of the sequence $u_n$ defined on a metric space admits a convergent sub-subsequence with limit $u$, then the original sequence $u_n$ converges to $u$ as well.}.		\end{proof} The proof of the lemma also shows also that the derivative or order $l-1$ is even locally Lipschitz in $B_{1,d} \setminus \{0\} $.


	










\begin{remark}
Taking $d=s$ in Theorem \ref{mainthm} guarantees the existence of a local change of coordinates around zero that conjugates the full flow map to its linear part in the presence of resonances, provided the specified regularity and spectral properties are satisfied. We further note that, thanks to the scaling in $\delta$ in the fixed-point argument, it is enough that $N\geq 2$, but for the resolution of all the resonances, i.e., the invertibility of $\mathcal{S}_A$, we have to go to higher orders in general. Of course, the order-by-order solution of the invariance equation allows us to go to arbitrarily large $N$. 
\end{remark}

We need another little lemma for the proof of the main theorem in the following section.




\begin{lemma}\label{lemma_coefficient_estimate}
Let $P(y)=\sum_{|q|\leq K}a_qy^q$ be a polynomial on $\mathbb C^M$. For every $\beta,\sigma>0$ there is a constant $C_0>1$, depending only on $M,\beta,$ and $\sigma$, such that
\begin{equation}
 |a_q|\leq C_0^K\|P\|_{K\beta,\sigma},\qquad |q|\leq K.
\end{equation}
\end{lemma}
\begin{proof}
Choose $0<r_0<\sigma$. The closed polydisc $\{|y_j|\leq r_0\}_{j=1}^M$ is contained in $\mathfrak T_\sigma$. On this polydisc, $|\Re y|\leq \sqrt M\,r_0$, and hence
\begin{equation}
 |P(y)|\leq e^{K\beta\sqrt M r_0}\|P\|_{K\beta,\sigma}.
\end{equation}
Cauchy's coefficient estimate on the polydisc gives
\begin{equation}
 |a_q|\leq r_0^{-|q|}e^{K\beta\sqrt M r_0}\|P\|_{K\beta,\sigma}
 \leq \bigl(e^{\beta\sqrt M r_0}\max\{1,r_0^{-1}\}\bigr)^K\|P\|_{K\beta,\sigma}.
\end{equation}
Taking $C_0=e^{\beta\sqrt M r_0}\max\{1,r_0^{-1}\}$ and enlarging it slightly if necessary proves the claim.
\end{proof}

\subsection{Proof of the Main Theorem as well as the stable Hartman conjecture}

In this section, we collect and organize the results from the previous section to complete the proof of the main theorem.  \\





\begin{proof}[Proof of Theorem~\ref{mainthm}]
Choose $N$ sufficiently large so that the logarithmic polynomial
$h^{<N}$ constructed in Proposition~\ref{approx}
contains all resonant indices and such that Proposition~\ref{Lemma_SM}
applies on the tail space $\Gamma_N^d$.
Here, as in Proposition~\ref{propfixedpoint}, $h^{<N}$ denotes the $\delta$-dependent approximate solution of the scaled equation and its homogeneous coefficient of order $k$ is $\delta^{k-1}\widehat h_k$. By Proposition~\ref{propfixedpoint}, after choosing the scaling parameter
$\delta>0$ sufficiently small, there exists a unique
$H\in\Gamma_N^d$ such that $h=h^{<N}+H$ solves the scaled invariance equation. It remains to establish the regularity of the correction $H$. By
construction, all resonant contributions of order smaller than $N$
are contained in $h^{<N}$. Hence, by
Lemma~\ref{lemma_regularity},
$h^{<N}\in C^{r,1-\varepsilon}$ for every
$\varepsilon\in(0,1)$, where $r=r_*$ is the uniform regularity index in Theorem~\ref{mainthm}. We claim that the tail $H$ has the same
regularity.\\
We first note that the fixed-point construction preserves the
polynomial-logarithmic structure used in Proposition~\ref{approx}.
Indeed, start the iteration by defining the fixed point in
Proposition~\ref{propfixedpoint} with $H_0=0$ and write
\begin{equation}
    H_{k+1}
    =
    T_\delta(H_k,h^{<N}).
\end{equation}
Since $h^{<N}$ is a logarithmic polynomial, analyticity of
$N_f^\delta$ implies that, at each fixed order $n$, the coefficient
of $x^n$ in $N_f^\delta(h^{<N}+H_k)$ is a polynomial in the
logarithmic variables and depends only on finitely many coefficients
of $h^{<N}$ and $H_k$. Moreover, the support condition of
Lemma~\ref{lemma_regularity} is preserved under products and sums.
Proposition~\ref{Lemma_SM} acts coefficient-wise and does not
introduce any additional logarithmic variables. Consequently, each
$H_k$ has an expansion
\begin{equation}
    H_k(x)
    =
    \sum_{|n|\geq N} H_{k,n}(\eta(x))x^n,
    \label{polynomial-tail-iterates}
\end{equation}
where the $H_{k,n}$ are polynomials satisfying the same support
condition as in Proposition~\ref{approx}.\\ 
More precisely, write
\begin{equation}
    H_{k,n}(\eta)
    =
    \sum_q a_{k,n,q}\eta^q .
\end{equation}
If $\eta_\ell$ was introduced at the resonant multi-index
$m^{(\ell)}$, then
\begin{equation}
    a_{k,n,q}\neq0
    \quad\Longrightarrow\quad
    n_j\geq
    \sum_{\ell=1}^M q_\ell m_j^{(\ell)},
    \qquad 1\leq j\leq d.
    \label{log-support-multiplicity}
\end{equation}
Indeed, this property holds for the approximate solution
$h^{<N}$ and is preserved under multiplication and addition.
Moreover, Proposition~\ref{Lemma_SM} only applies shifts in the
logarithmic variables and therefore can lower, but cannot increase,
the powers of the logarithmic variables. Hence
\eqref{log-support-multiplicity} is preserved by the fixed-point
iteration. Since $r\le r(m^{(\ell)})$ for every $\ell$, and the exponents $\nu_j$ were chosen in Proposition~\ref{approx} to satisfy the support and largeness conditions used in Lemma~\ref{lemma_regularity}, that lemma applies with regularity index $r$ to each factor $x^{m^{(\ell)}}\eta_\ell(x)$.
The order-by-order argument in the proof of
Proposition~\ref{approx} also gives
\begin{equation}
    \deg H_{k,n}\leq |n|-1.
    \label{tail-degree-bound}
\end{equation}
Since $H_k\to H$ in $\Gamma_N^d$, and for every fixed $n$ the polynomials satisfying \eqref{tail-degree-bound} form a finite-dimensional closed subspace of the corresponding coefficient space, the limiting coefficient $H_n$ is again a polynomial.
Therefore
\begin{equation}
    H(x)
    =
    \sum_{|n|\geq N} H_n(\eta(x))x^n,
    \qquad
    \deg H_n\leq |n|-1,
    \label{polynomial-tail}
\end{equation}
and the coefficients of $H_n$ satisfy
\eqref{log-support-multiplicity}.\\
We now prove that this series converges in $C^{r,\alpha}$ after possibly
restricting the domain. Fix $\alpha\in(0,1)$ and write
\begin{equation}
    H_n(\eta)
    =
    \sum_{|q|\leq |n|-1} a_{n,q}\eta^q.
\end{equation}
Applying Lemma~\ref{lemma_coefficient_estimate} with $K=|n|$, and noting that $\deg H_n\leq |n|-1$, gives a constant $C_0>1$, independent of $n$ and $q$, such that
\begin{equation}
    |a_{n,q}|
    \leq
    C_0^{|n|}
    \|H_n\|_{|n|\beta,\sigma}.
    \label{coefficient-bound}
\end{equation} For a coefficient $a_{n,q}\neq0$, set
\begin{equation}
    p_{n,q}
    :=
    n-\sum_{\ell=1}^M q_\ell m^{(\ell)}
    \in\mathbb{N}^d,
\end{equation}
which is well-defined by \eqref{log-support-multiplicity}. We may
then factor
\begin{equation}
    x^n\eta(x)^q
    =
    x^{p_{n,q}}
    \prod_{\ell=1}^M
    \left(
        x^{m^{(\ell)}}\eta_\ell(x)
    \right)^{q_\ell}.
    \label{tail-factorization}
\end{equation}
By Lemma~\ref{lemma_regularity},
\begin{equation}
    x^{m^{(\ell)}}\eta_\ell(x)
    \in C^{r,\alpha}
\end{equation}
for every $\ell$. Moreover, the estimates in the proof of
Lemma~\ref{lemma_regularity} show that
\begin{equation}
    \bigl\|
        x^{m^{(\ell)}}\eta_\ell(x)
    \bigr\|_{C^{r,\alpha}(B_\rho^d)}
    \longrightarrow0,
    \qquad
    \rho\longrightarrow0.
    \label{log-factor-small}
\end{equation}
Since $C^{r,\alpha}(B_\rho^d)$ is a Banach algebra, the
factorization \eqref{tail-factorization}, together with the
elementary estimates for the polynomial factor $x^{p_{n,q}}$,
implies that, after choosing $\rho>0$ sufficiently small, there
exist constants $C>0$ and $\vartheta\in(0,1)$, independent of
$n$ and $q$, such that
\begin{equation}
    \bigl\|
        x^n\eta(x)^q
    \bigr\|_{C^{r,\alpha}(B_\rho^d)}
    \leq
    C\vartheta^{|n|}.
    \label{monomial-tail-estimate}
\end{equation}
Indeed, taking absolute values in \eqref{tail-factorization} gives
\begin{equation}
    |n|
    =
    |p_{n,q}|
    +
    \sum_{\ell=1}^M q_\ell |m^{(\ell)}|.
\end{equation}
Let
\begin{equation}
    m_*:=\max_{1\leq \ell\leq M}|m^{(\ell)}|.
\end{equation}
Hence, for every pair $(n,q)$ with $a_{n,q}\neq0$, either
\begin{equation}
    |p_{n,q}|\geq \frac{|n|}{2},
\end{equation}
or
\begin{equation}
    \sum_{\ell=1}^M q_\ell |m^{(\ell)}|
    \geq \frac{|n|}{2},
\end{equation}
which implies
\begin{equation}
    |q|
    =
    \sum_{\ell=1}^M q_\ell
    \geq
    \frac{|n|}{2m_*}.
\end{equation}
In the first case, the polynomial factor $x^{p_{n,q}}$ yields
exponential decay in $|n|$ on a sufficiently small ball
$B_\rho^d$, while the polynomial losses in $|p_{n,q}|$ arising from
taking at most $r$ derivatives are absorbed by this exponential
decay after decreasing $\rho$ if necessary. In the second case,
\eqref{log-factor-small} and the Banach-algebra property of
$C^{r,\alpha}(B_\rho^d)$ give exponential decay through the factor
\begin{equation}
    \prod_{\ell=1}^M
    \bigl(x^{m^{(\ell)}}\eta_\ell(x)\bigr)^{q_\ell}.
\end{equation}
Consequently, after restricting $\rho>0$ further if necessary,
there exist constants $C>0$ and $\vartheta\in(0,1)$, independent
of $n$ and $q$, such that
\begin{equation}
    \bigl\|
        x^n\eta(x)^q
    \bigr\|_{C^{r,\alpha}(B_\rho^d)}
    \leq
    C\vartheta^{|n|}.
\end{equation}
Combining \eqref{coefficient-bound} and
\eqref{monomial-tail-estimate}, and decreasing $\rho$ further if
necessary so that the decay in \eqref{monomial-tail-estimate}
dominates the factor $C_0^{|n|}$ as well as the number of
monomials of degree at most $|n|-1$, we obtain a constant
$C_\rho>0$ such that
\begin{equation}
    \bigl\|
        H_n(\eta(\,\cdot\,))(\,\cdot\,)^n
    \bigr\|_{C^{r,\alpha}(B_\rho^d)}
    \leq
    C_\rho
    \|H_n\|_{|n|\beta,\sigma},
    \qquad |n|\geq N.
    \label{tail-regularity-estimate}
\end{equation}
Consequently,
\begin{align}
    \|H\|_{C^{r,\alpha}(B_\rho^d)}
    &\leq
    C_\rho
    \sum_{|n|\geq N}
        \|H_n\|_{|n|\beta,\sigma}
    \notag\\
    &=
    C_\rho\|H\|_{\Gamma_N^d}
    <\infty.
    \label{tail-holder-bound}
\end{align}
Hence $H\in C^{r,\alpha}(B_\rho^d)$. Since
$\alpha\in(0,1)$ was arbitrary, we conclude that
\begin{equation}
    H\in C^{r,1-},
    \qquad
    h=h^{<N}+H\in C^{r,1-}.
\end{equation}
Finally, $Dh(0)=\iota_{X_1}$, so the constant-rank theorem implies,
after restricting the domain if necessary, that $h$ is an embedding
and its image is an $f$-invariant manifold tangent to $X_1$ at the
origin. In the full-dimensional case $d=s$, $Dh(0)=I$, and the
inverse-function theorem yields a local conjugacy to the linearized
dynamics.
\end{proof}



Corollary~\ref{cor_Hartman} follows by taking $d=s$ in Theorem~\ref{mainthm} and applying the inverse-function theorem, as in the final paragraph of the proof.

%\textcolor{red}{At the first resonant index number $n_{\rm min}$, defined in \eqref{defNresnures}, Lemma \ref{shifted} applies and for each resonant index and we add $\Delta_{d-1,1} = d$ dimensions for each resonant index with index number $N_{\rm res}$.  This gives both property $(2)$ and $(3)$ in Theorem \ref{mainthm}. It follows that even if the map $f$ is analytic, the family of invariant manifolds guaranteed to exist by Theorem \ref{mainthm} will only be finitely-differentiable in the presence of resonances. }





\section{Examples and Counterexamples}\label{examples}

In this section, we present some examples and counterexamples that illustrate the assumptions and conclusions of our main theorems.

%\begin{example}	\begin{equation}	\begin{cases}	\dot{x}=-ax,\\	\dot{y}=-by,	\end{cases}	\end{equation}	for $a,b>0$, $x,y:\mathbb{R}\to\mathbb{R}$, with solution $\Big(x(t),y(t)\Big)=(x_0e^{-at},y_0e^{-bt})$. For $x_0\neq 0$, any trajectory runs on one of the invariant manifolds	\begin{equation}	y_q(x_0)=y_0\left(\frac{x}{x_0}\right)^q,	\end{equation}	depending only on the quotient $q=\frac{b}{a}$. If we are seeking a Taylor expansion of $y_q$ in $x_0$, i.e., assume that $y(t)=h(x(t))$, with $h(x)=h_2x^2+h_3x^3+\mathcal{O}(x^4)$, we obtain the equation	\begin{equation}\label{hn}	(na-b)h_n=0,	\end{equation}	for the $n^{th}$ coefficient $h_n$. In particular, in the absence of resonances, $h_n=0$ for all $n\geq 0$.\\	If, however, $q=N\in\mathbb{N}$, then equation \eqref{hn} has an additional solution and we can find a one-parameter family of solutions to the invariant equation of the form	\begin{equation}\label{yN}	y_N(x_0)=y_0x_0^{-N}x^N.	\end{equation}	Note that the invariant manifold $y_N$ in \eqref{yN} is also analytic and tangent to the $x$-axis.\\	The same reasoning can be applied if $q\in\mathbb{Q}$ and $y$ is expanded in a Puiseux series $y=h(x)=\sum_{n\geq n_0}h_nx^{\frac{n}{p}}$.\end{example}

\begin{example}
This example illustrates the simplest resonant system exhibiting logarithmic
terms. Consider the planar differential equation
\begin{equation}\label{explres}
\begin{cases}
\dot x=-x,\\
\dot y=-2y+x^2,
\end{cases}
\end{equation}
whose eigenvalues satisfy a $(1:2)$ resonance relation. The system can be solved explicitly. For the initial condition
$(x(0),y(0))=(x_0,y_0)$ one obtains
\begin{equation}\label{sol1}
\begin{split}
x(t)&=x_0e^{-t},\\
y(t)&=\bigl(y_0+x_0^2t\bigr)e^{-2t}.
\end{split}
\end{equation}
Eliminating the time variable via
\begin{equation}
t=-\log\!\left(\frac{x}{x_0}\right)
\end{equation}
gives the family of solution curves
\begin{equation}\label{soly1}
y=
\left[
y_0
-
x_0^2
\log\!\left(\frac{x}{x_0}\right)
\right]
\left(\frac{x}{x_0}\right)^2,
\end{equation}
or equivalently,
\begin{equation}
y
=
\frac{y_0}{x_0^2}x^2
-
x^2\log\!\left(\frac{x}{x_0}\right).
\end{equation}
Thus every trajectory with $x_0\neq0$ approaching the equilibrium, when represented as a graph over $x$, contains the resonant term
$x^2\log x$. The coefficient of the quadratic polynomial part depends on the
initial condition, whereas the logarithmic coefficient is fixed by the resonance. Differentiating twice shows that any non-trivial solution curve is necessarily of class
$C^{1,\alpha}$ for every $\alpha<1$, but not $C^2$ at the
equilibrium, in complete agreement with the regularity predicted by the main theorem.
\end{example}


\begin{example}
Let us illustrate the non-uniqueness of the resonant invariant manifolds in a little more detail. Consider the following two-dimensional dynamical system:
\begin{equation}\label{systexpl1}
\begin{cases}
\dot{x}=-x+y^2,\\
\dot{y}=-3y + a x^2+ b x^3,
\end{cases}
\end{equation}
for some parameters $a,b \in\mathbb{R}$, whose linear part has a $1:3$ resonance.\\
For $b\neq 0$, the Taylor coefficient of $x^3$ in the $y$-dynamics is non-zero and we can capture any local resonant invariant manifold only by expansion in logarithmic polynomials. For $b=0$, we will see that the logarithmic approximation reduces to an ordinary polynomial approximation. Since the eigenspace associated to the slow $x$-direction is just the $x$-axis, we can assume that the invariant manifold is a graph for $x$ small enough and we can thus expand $y$ as a logarithmic polynomial of the form
\begin{equation}\label{hexpl1}
\begin{split}
y(t)&=:h(x(t))\\
&=\sum_{n=1}^N h_n p_n(\log(x)) x^n(t)+\mathcal{O}(x(t)^{N+1}),
\end{split}
\end{equation}
for any $N\in \mathbb{N}$, $h_n\in \mathbb{R}$ and some real-valued polynomials $p_n$. Note that we only seek an expansion in the non-negative $x$-direction in \eqref{hexpl1}.\\
 Comparison of powers shows that
\begin{equation}\label{yexpl}
\begin{split}
h(x)&=a x^2  +  (-b \log(x) + \sigma)x^3  + a^3 x^5\\
&\quad+ \left(\frac{1}{9} a^2 (21 \sigma+4 b)-\frac{7}{3} a^2 b \log (x)\right) x^6\\
&\quad +\frac{a}{8}  \left(b^2+4 b \sigma +4 b \log (x) (4 b \log (x)-b-8 \sigma )+16
\sigma ^2\right)x^7\\
&\quad +\mathcal{O}(x^8),
\end{split}
\end{equation}
for any $\sigma\in\mathbb{R}$, defines an approximation to an invariant manifold. Note that the fifth-order term in the expansion \eqref{yexpl} does not contain any logarithmic terms, even though the third order term does. 
%In fact, the logarithmic terms only enter at multiplies of three. 
At zero, all the approximate invariant manifolds obtained above have the same tangency. If $b=0$, we see that, indeed, all logarithmic terms disappear. For $a=0$, we see that the higher-order terms, i.e., larger than order three, disappear. Figure \ref{plot1} illustrates the approximations \eqref{yexpl} for different values of $\sigma$ as well as the polynomial approximation up to order two.
\begin{figure}
	\includegraphics[scale=1.2]{res_manifold}\caption{Phase portrait of system \eqref{systexpl1} for parameter values $a=b=1$ together with the approximate invariant manifolds \eqref{yexpl} for $\sigma=0$ (red) and for $\sigma=1$ (green). The orange line shows the polynomial approximation $h(x)=x^2$ $(N=2)$. }
	\label{plot1}
\end{figure}
\end{example}


\begin{example}
    The following example illustrates that the $\alpha$-H\"{older} regularity guaranteed by the main theorem can, in general, not be improved. Consider the following three-dimensional differential equation,
    \begin{equation}\label{system_example}
        \begin{cases}
            \dot{x} & = - x,\\
            \dot{y} & = -2y,\\
            \dot{z} & = -3z + xy,
         \end{cases}
    \end{equation}
with explicit solution
\begin{equation}\label{solution}
\begin{split}
    x(t) & = x_0 e^{-t},\\
    y(t) & = y_0 e^{-2t},\\
    z(t) & = e^{-3t}(t x_0 y_0 + z_0). 
\end{split}
\end{equation}
Solving the first two equations in \eqref{solution} for the time gives 
\begin{equation}
    t = -\log\left(\frac{x}{x_0}\right) = -\frac{1}{2}\log\left(\frac{y}{y_0}\right),
\end{equation}
which implies the representation of time as one of the basis functions 
\begin{equation}\label{rep_time}
    -\frac{1}{2}\log(|x|^2+|y|)
    =t-\frac{1}{2}\log(|x_0|^2+|y_0|),
\end{equation}
The general form of an invariant manifold using the representation of time \eqref{rep_time}, conjugated to the $(x,y)$-dynamics, which are already linear in this example, is thus given by\begin{equation}
\label{equationz}
z = xy \Big(A -\frac{1}{2} \log(|x|^2+|y|) \Big),
\end{equation}
for any $A\in \mathbb{R}$. Figure \ref{plotz} shows the function \eqref{equationz} along with its derivatives. \\
Similarly, the flow map \eqref{solution} can be written as 
\begin{equation}
    (x,y,z)(t;x_0,y_0,z_0) = h(e^{-t}x_0,e^{-2t}y_0,e^{-3t}z_0),
\end{equation}
for the conjugacy map
\begin{equation}\label{h}
    h(x,y,z) = \left(\begin{array}{c}
        x\\ 
        y\\
        -\frac{1}{2}\log(|x^2|+|y|)xy + z
    \end{array}\right).
\end{equation}
Indeed, the functional form of \eqref{h} is guaranteed by our main theorem and the successive approximations in terms of logarithmic polynomials. In particular, we can read off the $C^{1,\alpha}$-regularity for any $\alpha\in(0,1)$ from \eqref{h} directly.\\
Let us show that the $C^{1,\alpha}$-regularity cannot be improved. More generally, a conjugacy map for system \eqref{system_example} of the form relevant to the above construction is given by
\begin{equation}
    h(x,y,z) = \left(\begin{array}{c}
        x\\ 
        y\\
        T(x,y) xy + z
    \end{array}\right),
\end{equation}
for a function $T:U\to\mathbb{R}$, defined on some open neighborhood $U\subset\mathbb{R}^2$ of the origin such that
 \begin{equation}
    T(x_0e^{-t},y_0e^{-2t}) = T(x_0,y_0)+t,\quad t\geq 0. 
\end{equation}
In particular, any such $T$ must be unbounded at the origin and thus the mixed $(x,y)$-derivative of the function $g(x,y) = xy T(x,y)$ cannot be bounded at the origin. Indeed, 
\begin{equation}
    g(e^{-t}x,e^{-2t}y)
=
e^{-3t}\bigl(g(x,y)+txy\bigr).
\end{equation}
Assume, for contradiction, that $g$ extends to a $C^2$ function in a neighborhood of the origin. Differentiating the preceding identity once with respect to $x$ and once with respect to $y$ gives
\begin{equation}
    e^{-3t}g_{xy}(e^{-t}x,e^{-2t}y)
=
e^{-3t}\bigl(g_{xy}(x,y)+t\bigr),
\end{equation}
and hence
\begin{equation}
g_{xy}(e^{-t}x,e^{-2t}y)
=
g_{xy}(x,y)+t.
\end{equation}
For any fixed $(x,y)$ with $xy\neq0$, the point $(e^{-t}x,e^{-2t}y)$ converges to the origin as $t\to\infty$, whereas the right-hand side diverges linearly in $t$.
 This contradicts the boundedness of $g_{xy}$ near the origin. Therefore, for every solution $T$ of the cohomological equation, the function $xyT(x,y)$ fails to extend as a $C^2$ function at the origin.





%Clearly, for any $(x_0,y_0,z_0)\in\mathbb{R}^3$, we can find an $A\in\mathbb{R}$ (choosing $B=0$) such that \eqref{equationz} holds, thus showing that the parametrization \eqref{equationz} indeed exhausts the whole phase space. 
\begin{figure}
    \centering
    \includegraphics[width=0.45\linewidth]{pf.png}
    \includegraphics[width=0.45\linewidth]{pfx.png}
    \includegraphics[width=0.45\linewidth]{pfy.png}
    \includegraphics[width=0.45\linewidth]{pfxy.png}
    \caption{The invariant manifold \eqref{equationz} (blue) with its first derivatives (orange) and its $xy$-derivative (red), showing that the resonant manifold is, indeed, $C^{1,\alpha}$, but not $C^2$. }
    \label{plotz}
\end{figure}
\end{example}









%\begin{example} We have a $1:3$ resonance, but the right-hand side is only $C^2$.\end{example}



\section{Conclusion and Further Perspectives}\label{conclusion}


We have proved the existence and regularity of stable invariant manifolds and linearization maps for analytic dynamical systems near a hyperbolic fixed point in the presence of resonances, including external resonances.  The main conclusion is that external resonances need not be viewed only as obstructions to the invariance equation.  Rather, they indicate that the usual Taylor class is too restrictive.  Once ordinary polynomial expansions are replaced by logarithmic-polynomial expansions, the resonant cohomological equations can be solved recursively to arbitrary finite order and the resulting approximate solution can be completed by a fixed-point argument.\\

This gives a precise description of the regularity loss caused by resonance.  In the non-resonant case one recovers uniqueness and analyticity of the invariant manifold \cite{CABRE2005444}.  In the externally resonant case one generally obtains non-unique invariant manifolds and conjugacy maps of finite differentiability, even for analytic maps. The resonant indices determine a uniform spectrum-based regularity guarantee, yielding $C^{r,1-\varepsilon}$ regularity for every $\varepsilon\in(0,1)$. When the corresponding resonant forcing coefficient is nonzero, the examples show that the resulting logarithmic loss can be intrinsic, while if that coefficient vanishes, higher regularity may occur.  Logarithmic terms are therefore the natural local signature of the external resonance mechanism.\\

The proof separates the algebraic and analytic aspects of the problem.  The algebraic step consists of solving the resonant cohomology equations order by order in a logarithmic-polynomial algebra.  The analytic step uses the resulting approximate solution as the starting point for a contraction argument in a suitable H\"{o}lder space.  This separation makes the method explicit and, in principle, algorithmic.  It suggests computational procedures for resonant invariant manifolds, resonant spectral submanifolds, and local conjugacy maps beyond the non-resonant regime.\\

Several extensions are natural.  First, one should remove some simplifying assumptions of the present work.  Although we have formulated the results in the analytic category, the resonance mechanism itself is not essentially analytic, and analogous statements should hold under sufficiently high finite smoothness.  Similarly, the assumption that the linear part is semisimple should be relaxed.  For hyperbolic matrices with Jordan blocks, the nilpotent part of the linear dynamics is expected to generate additional polynomial factors, which should interact with the logarithmic terms produced by resonances.  A complete treatment of this case would give a more general local normal form for resonant invariant manifolds.\\

Another direction is the extension from the stable setting to mixed hyperbolic spectra and to the full Hartman conjecture as well as the Van Strien conjecture \cite{vanstrien1990smooth}.  Stable and unstable spectra are formally related by inversion, but a simultaneous treatment of stable and unstable directions is more delicate \cite{hartman1960lemma,hartman1960local,newhouse2017hartman,pugh1969theorem}.  In particular, internal and external resonances may interact with both forward and backward dynamics, and some resonant terms may be absorbed into the reduced dynamics while others remain transverse obstructions.  A logarithmic-polynomial version of local conjugacy theory could provide a new approach to optimal regularity questions for Hartman--Grobman conjugacies and invariant foliations \cite{bates2000invariant,hasselblatt1994periodic,pugh1997holder,luZhangZhang2017differentiability}.\\

The construction also suggests applications in areas where resonant invariant manifolds play an organizing role.  In celestial mechanics, stable and unstable manifolds of libration-point dynamics are central in the design of low-energy transfers and in the study of resonance transitions in the restricted three-body problem \cite{koon2000heteroclinic,koon2001low,koon2011dynamical,gomez2004connecting}.  Logarithmic-polynomial parameterizations may provide more accurate local descriptions near resonant libration-point orbits and normally hyperbolic structures.  In chemical reaction dynamics, normally hyperbolic invariant manifolds and their stable and unstable manifolds organize phase-space transport and transition-state theory \cite{wiggins1994normally,komatsuzaki2001dynamical,waalkens2004direct,ezra2009microcanonical}. Since resonances affect the local normal forms near saddle-type equilibria, the present framework may help describe the asymptotic structure of resonant transition-state manifolds.\\

Another important application is nonlinear model reduction for mechanical systems.  Spectral submanifolds provide a rigorous invariant-manifold interpretation of nonlinear normal modes and an exact reduced dynamics near equilibria or periodic orbits \cite{shaw1991normal,shaw1993normal,haller2016nonlinear}.  Existing SSM methods rely heavily on polynomial parameterizations and on non-resonance or spectral-quotient conditions \cite{szalai2017nonlinear,breunung2018explicit,jain2018exact,ponsioen2020model,cenedese2022data}.  The present results suggest a resonant extension of this framework in which external resonances are handled by enlarging the parameterization class to include logarithmic terms.  This could lead to effective reduced-order models in regimes where classical polynomial SSM expansions fail or lose differentiability.\\

Finally, it would be interesting to extend the theory to functional differential equations and other infinite-dimensional systems.  Delay equations possess finite-dimensional invariant manifolds associated with spectral subspaces of the linearized semi-flow, and resonances already play a central role in their normal-form theory \cite{hale1993introduction,dieckmann1995delay,faria1995normal}.  At the same time, the Banach-space setting contains genuinely infinite-dimensional obstructions: even for contractions, $C^1$ linearization can fail without additional assumptions \cite{rodrigues2004linearization,rodrigues2005invertible,rodrigues2012known}.  A useful long-term goal is therefore to distinguish finite-dimensional resonance effects, which may be resolved by logarithmic-polynomial expansions, from obstructions caused by the spectrum and geometry of the ambient infinite-dimensional space.\\

Overall, the results show that external resonances identify the asymptotic structure that replaces analyticity.  Once this structure is incorporated into the functional setting, existence and optimal regularity become accessible, opening a path toward a broader theory of resonant invariant manifolds with applications to local linearization, celestial mechanics, chemical dynamics, functional differential equations, and nonlinear model reduction.\\

\section{Acknowledgments}
The authors declare no conflict of interest. 






\bibliographystyle{abbrv}
\bibliography{DynamicalSystems}


\appendix

\section{Relation To Flow Maps of Differential Equations}\label{sec_appendix}
As we prefer to formulate our results in terms of maps, we will briefly recall the relation to differential systems in this section, which corresponds to studying the flow map for a conveniently chosen time $T$. 
Consider the dynamical system
\begin{equation}\label{dynsyst}
\dot{x}=\mathcal{X}(x), 
\end{equation}
defined on an open subset $U\subseteq\mathbb{R}^s$ containing the origin for an analytic vector field $\mathcal{X}$ such that $\mathcal{X}(0)=0$. We write
\begin{equation}\label{dyn}
\mathcal{X}(x)=\mathcal{B}x+\mathcal{N}(x),
\end{equation}
where $\mathcal{B}=D\mathcal{X}(0)$ and $\mathcal{N}(x)=\mathcal{O}(|x|^2)$. For each $t\geq 0$, equation \eqref{dyn} defines an analytic flow map $F^t:U\mapsto \mathbb{R}^s$ such that
\begin{equation}\label{time1} 
 F^t(0) = 0; \quad DF^t(0)=e^{\mathcal{B}t},
\end{equation}
i.e., here we have that $A=e^{\mathcal{B}t}$ in the notation of \eqref{defA}.\\
We will show in the following that the results on maps imply the results for flows under the same hypotheses. We note, however, that maps appear in other applications besides the time-one maps of a differential equation as well, e.g., as return maps to a surface of section. Because of \eqref{time1}, a map with $\det DF(0) < 0$ cannot appear as a time-one map of a differential equation, which implies that the results for maps are, in principle, more general.\\
The main hypotheses in the theorem are twofold: regularity and resonances. Clearly, the regularity of the flows implies the regularity of the time-$T$-maps. As for the resonances, we recall that if  the eigenvalues of $\mathcal{B}$ are $\mu_1, \ldots, \mu_D$, the eigenvalues of $\exp( T \mathcal{B})$ are $\lambda_1 = \exp(T \mu_1), \ldots, \lambda_D = \exp(T\mu_D)$. Denoting $\mu=(\mu_1,..,\mu_D)$, a resonance for the vector field occurs whenever
\begin{equation} \label{flowresonance} 
\mu \cdot n  = \tilde \mu,  
\end{equation}
for some $n \in \mathbb{N}^D$ with $|n|\geq2$ and some $\tilde \mu\in \{\mu_1,...\mu_D\}$. On the other hand, a resonance for the time $T$ map occurs when 
\begin{equation} \label{mapresonance} 
\lambda^n = \tilde \lambda,
\end{equation}
for some $n \in \mathbb{N}^D$ with $|n|\geq2$ and some $\tilde \lambda \in \{\lambda_1,...,\lambda_D\}$.\\
Since $\lambda^n =  \exp(T \mu \cdot n)$, a resonance for the flow implies a resonance for the time-$T$-map for any $T\geq 0$.  As for the converse, we note that a resonance for the time-$T$-map does not necessarily imply the existence of a resonance for the flow \eqref{flowresonance}. Only the weaker equation $ T \mu \cdot n  =  T\tilde \mu + 2 \pi \ri k$ for some $k \in \mathbb{Z}$ holds. 
Hence, for a fixed $T$, the time-$T$ map may contain additional resonances that are not resonances of the vector field. These additional relations arise only when
\begin{equation}
T(\mu\cdot n-\widetilde\mu)\in 2\pi\ri\mathbb{Z}\setminus\{0\}.
\end{equation}
For a fixed finite collection of multi-indices, the exceptional values of $T$ satisfying one of these relations form a discrete set. We may therefore choose $T>0$ outside this set, so that, up to the finite order relevant for the construction, the resonances of the time-$T$ map are exactly those inherited from the flow. \\

\begin{remark}
Let $F^t$ denote the local flow generated by an analytic vector field with linearization $\mathcal{B}$ at the origin, so that $DF^t(0)=e^{t\mathcal{B}}$. Suppose that, for some $T>0$, a local $C^{r,\alpha}$-diffeomorphism $h$, normalized to $Dh(0)=I$, conjugates the time-$T$ map to its linearization,
\begin{equation}
h\circ F^T=e^{T\mathcal B}\circ h.
\end{equation}
Then this conjugacy can be promoted to a conjugacy of the full local flow. Indeed, define
\begin{equation}
\widetilde h(x)
:=
\frac{1}{T}\int_0^T e^{-t\mathcal{B}}h(F^t(x))\,dt.
\end{equation}
The map
\begin{equation}
t\mapsto e^{-t\mathcal B}h(F^t(x))
\end{equation}
is $T$-periodic by the time-$T$ conjugacy. Hence, for all $s$ for which the local flow is defined,
\begin{equation}
\widetilde h(F^s(x))
=
\frac{e^{s\mathcal{B}}}{T}
\int_s^{T+s} e^{-t\mathcal{B}}h(F^t(x))\,dt
=
e^{s\mathcal{B}}\widetilde h(x).
\end{equation}
Moreover,
\begin{equation}
D\widetilde h(0)
=
\frac{1}{T}\int_0^T
e^{-t\mathcal{B}}Dh(0)DF^t(0)\,dt
=
I.
\end{equation}
Therefore, after restricting the domain if necessary, $\widetilde h$ is a local $C^{r,\alpha}$-diffeomorphism. Thus a $C^{r,\alpha}$ conjugacy of a time-$T$ map yields a $C^{r,\alpha}$ conjugacy of the full local flow.
\end{remark}


\medskip
\noindent
%Let us also note that if there exists a \emph{unique} invariant manifold $M$ for the flow map $f^T$ for a specific time $T$ under some assumptions, the manifold $M$ is actually invariant for all times. Indeed, thanks to the flow property, we have that\begin{equation}f^T(f^t(M))=f^t(f^T(M))\end{equation}which shows that $f^t(M)$ is an invariant manifold of $f^T$ as well. Uniqueness of $M$ then implies that $f^t(M)\subseteq M$ for all $t$ sufficiently small. 



 



\end{document}
